% !TeX program = lualatex
% !TeX program = lualatex
\documentclass[11pt,a4paper]{ltjsarticle}
\usepackage[margin=25mm,headheight=17pt]{geometry}
\usepackage[haranoaji]{luatexja-preset}
\usepackage{amsmath,amssymb,amsthm,mathtools,mathrsfs,booktabs,array,enumitem,longtable}
\usepackage[unicode,hidelinks]{hyperref}
\usepackage{fancyhdr}
\pagestyle{fancy}\fancyhf{}\fancyfoot[C]{\thepage}
\fancyhead[L]{\small\nouppercase{\leftmark}}
\renewcommand{\sectionmark}[1]{\markboth{#1}{}}
\setlength{\emergencystretch}{3em}\allowdisplaybreaks[2]
\setlength{\parskip}{.15em}
\newtheoremstyle{jp}{.6em}{.6em}{\normalfont}{}{\bfseries}{.}{.5em}{}
\theoremstyle{jp}
\newtheorem{theorem}{定理}[section]
\newtheorem{proposition}[theorem]{命題}
\newtheorem{lemma}[theorem]{補題}
\newtheorem{corollary}[theorem]{系}
\newtheorem{definition}[theorem]{定義}
\newtheorem{remark}[theorem]{注意}
\newtheorem{example}[theorem]{例}
\newtheorem{question}[theorem]{問題}
\renewcommand{\proofname}{証明}
\newcommand{\kk}{\mathbb C}\newcommand{\PP}{\mathbb P}
\newcommand{\OO}{\mathcal O}\newcommand{\A}{\mathcal A}
\newcommand{\B}{B_{\mathrm{VK}}}\newcommand{\Blo}{B_{p,q}}
\newcommand{\HH}{\operatorname{HH}}\newcommand{\Hom}{\operatorname{Hom}}
\newcommand{\End}{\operatorname{End}}\newcommand{\Sym}{\operatorname{Sym}}
\newcommand{\Perf}{\operatorname{Perf}}\newcommand{\Der}{\operatorname{Der}}
\newcommand{\Inn}{\operatorname{Inn}}\newcommand{\gr}{\operatorname{gr}}
\newcommand{\rk}{\operatorname{rank}}\newcommand{\Spec}{\operatorname{Spec}}
\newcommand{\id}{\operatorname{id}}\newcommand{\Ext}{\operatorname{Ext}}
\newcommand{\Tor}{\operatorname{Tor}}\newcommand{\Ann}{\operatorname{Ann}}
\newcommand{\coker}{\operatorname{coker}}\newcommand{\Supp}{\operatorname{Supp}}
\newcommand{\pd}{\operatorname{pd}}\newcommand{\depth}{\operatorname{depth}}
\newcommand{\Frac}{\operatorname{Frac}}\newcommand{\Aut}{\operatorname{Aut}}
\newcommand{\Out}{\operatorname{Out}}\newcommand{\GL}{\operatorname{GL}}
\newcommand{\PGL}{\operatorname{PGL}}\newcommand{\Hilb}{\operatorname{Hilb}}
\newcommand{\dashsim}{\mathrel{\overset{\mathrm{bir}}{\sim}}}
\setcounter{tocdepth}{1}

\newcommand{\Fitt}{\operatorname{Fitt}}
\newcommand{\Nm}{\operatorname{Nm}}
\title{有限点族の相対 Serre 双対性と corank 二の幾何\\
\large 分岐自由因子、critical族内の非critical被覆と循環障害}
\author{}\date{2026年10月29日\quad 日本語原稿 第25版}
\begin{document}
\maketitle
\begin{abstract}
四頂点の Veronese--Kronecker 型代数の十六変数変形に付随する、
長さ十一の有限平坦点族 $T=B[[x,y]]/(F_1,F_2)$ を調べる。
普遍 thin module の相対三次 Ext が全族上で自由階数一であることから、
その Serre 固定性を閉点から全族へ持ち上げる。
自己 Ext 複体を二次正方 Jacobian へ縮約し、その鎖自己双対性を利用して
$F_{2y}\in(F_1,F_2,F_{1x},F_{1y},F_{2x})$ を全次数で証明する。
さらに同じ自己双対性を平均化し、点族全体を相対 almost closed 一形式の零スキームとして実現する。
従って corank 二の環は、十三変数環 $A_0$ 上自由階数四の正則局所環となる。
この四枚被覆は、長さ四の完全交叉
$\mathbb C[[X,Y]]/(XY,Y^2-2X^2)$ の四変数最小 versal 変形に
九変数の形式円板を加えたものであり、一般 Galois 群は $S_4$、
ramification divisor は三次元普通二重点と九次元形式円板の積である。
基底のbranch divisorは明示的な四次行列式で定まる被約既約自由因子であり、
その正規化はこの普通二重点である。
さらに四枚被覆自体には相対対称perfect obstruction theoryが存在しないことを証明し、
$(w,z)$ の一次元方向に沿う三変数sliceを分類する。
二つの例外方向を除く全てのsliceには、$D_4$ 型の三次式を中心とする
最小versalな多項式potentialを明示できる。
さらに、明示的なpotentialを持つ別の長さ十一の有限平坦族を構成し、
そのcorank二被覆が同じ非criticalな四枚被覆になることを示す。
従ってこの被覆の非critical性から、元の十一点族のpotential非存在は導けない。
また $Q=\Omega^1_{T/B}/dT$ と全 potential 障害加群 $M$ は、
$A_0$ 上で自由階数四の部分と循環 torsion に分解する。
正規化した相対 potential の存在は、この循環部分における
一つの元の一つのイデアルへの所属に帰着するが、その所属自体は未決定である。
二つの循環torsionイデアルの高さは少なくとも四である。
一般論として、標数零の有限可換代数族に対する $Q$ の台が
非 curvilinear ファイバーの軌跡の像と一致することも示す。
\end{abstract}
\tableofcontents

\section{序論}

有限点族の中心が potential の critical locus であるとき、
その potential を基底全体へ延長できるかという問題には、
点族の相対微分と、点対象の導来自己準同型の両方が関わる。
本稿ではこの二つを結び付け、十六変数の実際の変形族に対して
Jacobian の四つの成分のうち一つが他の三つに従属することを証明する。
この従属は全形式の恒等式であり、高次 jet の消滅の外挿ではない。

出発点は、Loginov--Orlov の非可換 Fano 三次元多様体の構成
\cite{LO}と、四頂点代数の関係変形との比較である。
本稿で固定する代数 $A_*$ は、滑らかな十六次元 miniversal 基底を持つ。
その平面点配置の一つの非被約部分は長さ十一であり、
対応する有限平坦族を
\[
 B=\kk[[t_1,\ldots,t_{16}]],\qquad
 T=B[[x,y]]/(F_1,F_2)
\]
と書く。中心の最低次関係は $x^2y,xy^2$ である。
この具体的な族では、関係、syzygy、普遍 thin module が同時に変形するため、
二つの任意の形式級数では得られない自己双対性を持つ。

主結果の証明は三つの独立した段階からなる。
まず syzygy-dual incidence の恒等式と単位小行列により、
普遍 thin module の三次自己 Ext が全点族上で自由階数一となることを示す。
次に相対 Serre 双対性から、閉点の Serre 固定性を全点族へ持ち上げる。
最後に自己 Ext 複体の実際の射影分解を縮約し、
可逆行列 $U,V$ による
\[
 UJ_F=J_F^tV
\]
を得る。中心の二次項を比較すると $U_{12}$ が単元であることが分かり、
\[
 F_{2y}\in(F_1,F_2,F_{1x},F_{1y},F_{2x})
\]
が従う（定理\ref{thm:psi-vanishing}）。
この方法は全環上で $U=V$ を要求せず、相対 potential の存在を仮定しない。
さらに $U$ と $V^t$ の平均を取ると、Jacobian を対称化する可逆行列を得る。
これにより点族は全形式の相対 almost closed 一形式の零スキームとして実現される
（命題\ref{prop:relative-almost-closed}）。

相対微分余核
\[
 Q=\Omega^1_{T/B}/dT
\]
は $B$-加群であり、相対 de Rham 第一コホモロジーとは区別する。
一般の有限可換代数族について、その台はファイバーの非 curvilinear 軌跡の像に一致する
（定理\ref{thm:support}）。二変数の場合には
$Q\twoheadrightarrow\Omega^2_{T/B}$ があり、その台は Jacobian の corank 二の軌跡である。
本例では
\[
 A_0=\kk[[t_i:i\notin\{5,15,16\}]],\qquad
 \Omega^2_{T/B}\simeq\Lambda
 =B[[x,y]]/(F_1,F_2,F_{1x},F_{1y},F_{2x})
\]
となり、$\Lambda$ は十三次元正則局所環、$A_0$ 上有限自由階数四である。
その Kodaira--Spencer 写像は全射であり、四枚被覆の全形式標準形、
普通二重点の ramification divisor、一般 Galois 群 $S_4$ が得られる
（第\ref{sec:four-versal-v23}節）。
続いて基底内の分岐因子をtrace判別式として具体化し、
その対数的ベクトル場の基底と、普通二重点による正規化を与える
（第\ref{sec:four-branch-v24}節）。
この自由性はversal有限写像の一般論にも従うが、本稿では全ての行列とnormを明示する。

四枚被覆自身のcritical性は、元の十一点族とは異なる結論になる。
相対余接複体の最小二項表示に対する一次係数の比較により、
四変数の全族は相対対称perfect obstruction theoryを持たないことを証明する。
この非存在はambientの次元を増やしても解消しない。
一方、$(w,z)=(pt,qt)$ という三変数sliceでは、critical性は
$p^2+2q^2\ne0$ と同値であり、その場合のpotentialを全次数の多項式で与える
（第\ref{sec:four-noncritical-v24}節）。

この相違を明示例でも示す。七変数基底上に全potentialを持つ長さ十一の有限平坦族を構成すると、
そのcorank二の軌跡は四変数最小versalな非critical四枚被覆になる。
九個の自明な形式変数を加えれば、基底の次元も実際の族の十六、十三と一致する
（命題\ref{prop:critical-source-noncritical-corank-v25}）。
したがってcorank二被覆だけの非critical性は、元の族のpotential非存在を結論する障害にはならない。

全相対 potential の問題には、別の有限加群が残る。
相対微分を十九次の単位小行列で消去すると、$3\times12$ 行列の余核 $M$ と
元 $\beta\in M$ が得られ、存在条件は $\beta=0$ となる。
本稿はこの問題をさらに
\[
 M\simeq A_0^4\oplus A_0/J_M,\qquad
 Q\simeq A_0^4\oplus A_0/(J_M,b_0),\qquad
 \beta\longleftrightarrow[b_0]
\]
へ帰着する。$J_M$ は三つの独立な線形項を持ち、$b_0$ は二次以上に取れる。
さらに最小関係数を二つの基底について比較すると、$J_M$ の最小生成数は少なくとも四となる。
三つの一次関係を消去した後にも高次関係が必ず残り、$J_M,J_Q$ の高さは少なくとも四である。
従って射影次元は $\pd_BM,\pd_BQ\ge7$ である。
これらの分解は $A_0$-線形なものであり、$B$-線形な分解を主張するものではない。

有限射の ramification module
$\mathscr R_{T/B}=\ker(d:T\to\Omega^1_{T/B})$ と opening の理論は
Ishikawa により研究されている \cite{IshikawaOpenings,IshikawaFrontal}。
本稿の存在問題も、指定した中心 class を持つ $\mathscr R_{T/B}$ の元を求める問題である。
新しい点は ramification module という概念の導入ではなく、
この非可換代数の変形から生じる有限点族における
Serre 双対性、正則な corank 二の環、四枚被覆の幾何、循環障害の具体的な連結にある。
全相対 potential の存在、すなわち $b_0\in J_M$ は本稿では決定しない。

\section{有限族の微分余核と非 curvilinear 軌跡}\label{sec:support}

\subsection{基底変更とファイバーによる判定}

$B$ を標数零の体上の Noether 環、$T$ を有限可換 $B$-代数とする。
以下で $Q$ は $B$-加群であり、通常は $T$-加群ではない。
\[
Q_{T/B}=\coker\bigl(d:T\longrightarrow\Omega^1_{T/B}\bigr).
\]
有限 $B$-加群なので、その台をファイバーで判定できる。
ここで有限型局所零次元代数が \emph{curvilinear} であるとは、
代数閉包への係数拡大後、各局所因子の埋込次元が高々一であることとする。

\begin{theorem}[微分余核の台]\label{thm:support}
$\pi:\Spec T\to\Spec B$ と置くと
\[
\Supp_B Q_{T/B}
=\pi\!\left(V\bigl(\Fitt_1^T\Omega^1_{T/B}\bigr)\right).
\]
すなわち $Q_{T/B}$ が零でない素点は、幾何学的ファイバーに
埋込次元二以上の局所因子が現れる素点と正確に一致する。
この主張には $T/B$ の平坦性を仮定しない。
\end{theorem}

\begin{proof}
任意の素イデアル $\mathfrak p\subset B$ について、微分の基底変更と余核の右完全性から
\[
Q_{T/B}\otimes_B\kappa(\mathfrak p)
\simeq
\Omega^1_{A/\kappa(\mathfrak p)}/dA,
\qquad A=T\otimes_B\kappa(\mathfrak p).
\]
さらに代数閉包への忠実平坦な係数拡大を行い、$A$ を局所 Artin 因子へ分解する。
微分とその余核はこの積分解に沿って分解する。

そこで $A$ を代数閉体 $K$ 上の局所 Artin 代数、極大イデアルを $\mathfrak m$ とする。
埋込次元が零なら $A=K$ である。埋込次元が一なら
$A\simeq K[z]/(z^e)$ であり、$e\ge2$ の場合
\[
\Omega^1_{A/K}=K[z]/(z^{e-1})\,dz.
\]
標数零なので $z^i\,dz=d(z^{i+1}/(i+1))$ であり、$dA$ が全体を生成する。

逆に $\dim_K\mathfrak m/\mathfrak m^2\ge2$ とする。
独立な二元を代表する $x,y\in\mathfrak m$ を取れば、
$dx\wedge dy$ の $\Omega^2_{A/K}\otimes_A K$ における像は非零である。
$x\,dy$ が完全形式なら、その外微分 $dx\wedge dy$ は零となるはずなので矛盾する。
従って $\Omega^1_{A/K}/dA\ne0$ である。

最後に、有限表示加群 $\Omega^1_{T/B}$ が剰余体上で二次元以上となる点は
$V(\Fitt_1^T\Omega^1_{T/B})$ である。
有限射の像は閉であり、有限加群に対する Nakayama の補題による台の判定を用いればよい。
\end{proof}

\subsection{二変数表示における外微分}

以下では $B$ は標数零の完備局所環とし、
\[
S=B[[x,y]],\qquad I=(F_1,F_2),\qquad T=S/I
\]
が $B$ 上有限であるとする。相対微分は連続微分を用いる。
有限 $B$-代数に移った後には通常の有限表示の相対微分と一致する。
相対 Jacobian を
\[
J_F=\begin{pmatrix}F_{1x}&F_{1y}\\F_{2x}&F_{2y}\end{pmatrix},
\qquad
\mathfrak j=(F_{1x},F_{1y},F_{2x},F_{2y})T
\]
と置く。

\begin{proposition}\label{prop:omega2}
$B$-線形な完全列
\[
0\longrightarrow H^1(\Omega^\bullet_{T/B})
\longrightarrow Q_{T/B}\xrightarrow{\ d\ }
T/\mathfrak j\longrightarrow0
\]
が存在する。最後の同型には $dx\wedge dy$ を用いる。
また
\[
\Supp_BQ_{T/B}=\Supp_B(T/\mathfrak j).
\]
\end{proposition}

\begin{proof}
$\Omega^1_{T/B}$ の表示を外積すると
$\Omega^2_{T/B}\simeq(T/\mathfrak j)\,dx\wedge dy$ を得る。
任意の $a\in T$ を $\widetilde a\in B[[x,y]]$ へ持ち上げる。
標数零なので $\partial_xG=\widetilde a$ を満たす形式級数 $G$ が存在し、
$d(G\,dy)=a\,dx\wedge dy$ である。
従って $d:\Omega^1_{T/B}\to\Omega^2_{T/B}$ は全射であり、完全形式を殺すので
表示された完全列が得られる。
台の等号は定理\ref{thm:support}と
$\Fitt_1^T\Omega^1_{T/B}=\mathfrak j$ による。
\end{proof}

これは台の等号であって、$Q_{T/B}=T/\mathfrak j$ という加群の同型ではない。
後述する具体例では、その核が非零の循環 torsion 加群であることまで分かる。

\section{具体的な点族と有限入力}\label{sec:input}

\subsection{四頂点代数と relation slice}

四頂点 $0,1,2,3$、隣接矢空間の次元が $(3,6,3)$ の箙を取る。
矢をそれぞれ $x_1,x_2,x_3$、
$u_{11},u_{12},u_{13},u_{22},u_{23},u_{33}$、$z_1,z_2,z_3$ と書く。
付録\ref{app:relations}の七つずつの左・右二次関係で定まる代数を $A_*$ とする。
この代数は $\HH^\bullet(A_*)=(1,0,16,21)$ を持ち、
miniversal 代数変形の基底は
\[
 B=\kk[[t_1,\ldots,t_{16}]],\qquad\mathfrak n=(t_1,\ldots,t_{16})
\]
である。以後の座標は、次の relation slice で固定する。

左右の七関係の Grassmann chart は百五十四座標を持つ。
三次の重なりを記述する $54\times42$ 行列を $N(v)$ とし、
中心での六次元核の座標を固定する。
階数三十六の minor を用いて核を解いた後、残る百八個の incidence 方程式の
微分は階数八十七である。
$\HH^2$ の十六代表を $H$、八十七個の pivot 方向を $E$ と書き、
\[
 v=Ht+Ew(t),\qquad w(0)=0
\]
とする。指定した八十七方程式を形式陰関数定理で解き $w(t)$ を取る。
関係スキームの $A_*$ における滑らかさから、残りの incidence 方程式も成立する。
座標と選択した minor は
\nolinkurl{archive/v21/code/full_relation_cubic_jets.json} および
\nolinkurl{archive/v21/code/relative_algebraization_frame_v19_certificate.json}
で指定される。従って本稿でいう「実際の族」は、この全次数の構成を指す。

\subsection{二つの determinantal 関係}

薄い点加群の chart $x_3=u_{33}=z_3=1$ において $x=x_1,y=x_2$ と置く。
左七関係を $\mathsf L(x,y)u=0$ と書くと、$\mathsf L$ は七行六列行列である。
行 $0,2,4,5,6$ と列 $0,1,2,3,4$ の小行列を $C$ と書く。
その閉点での determinant は一なので、五つの $u$ を消去できる。
行の順序を保持して
\[
 F_1=-\det\mathsf L_{(0,2,4,5,6,3),\,(0,1,2,3,4,5)},\qquad
 F_2=-\det\mathsf L_{(0,2,4,5,6,1),\,(0,1,2,3,4,5)}
\]
と定義する。これは消去後の残る左二関係の共通単元 $-\det C$ 倍である。
右の行 $0,2$ における $z_1,z_2$ の係数の $2\times2$ determinant も閉点で一であり、
右座標を一意に復元できる。
中心 $(t,x,y)=0$ を含む局所因子の完成を
\[
 S=B[[x,y]],\qquad I=(F_1,F_2),\qquad T=S/I
\]
と書く。係数 $w(t)$ は代数的形式級数であり、$F_i$ は $x,y$ の多項式である。
中心の最低次項は
\[
 F_1(0,x,y)=x^2y+O((x,y)^4),\qquad
 F_2(0,x,y)=xy^2+O((x,y)^4)
\]
となる。

中心のイデアルは、例えば次の七次多項式の Jacobian イデアルである。
\begin{align*}
W_7={}&\tfrac12x^2y^2+x^4y+xy^4-\tfrac15(x^5+y^5)\\
&+\tfrac12(x^6+y^6)-x^5y-xy^5-x^4y^2+11x^3y^3-\tfrac53x^2y^4\\
&-\tfrac37x^7-\tfrac{17}{3}x^6y+\tfrac{55}{3}x^5y^2
-\tfrac23x^4y^3-2x^3y^4+19x^2y^5-9xy^6+\tfrac17y^7.
\end{align*}
全相対問題で用いる正規化中心 class は
\[
 W_9=W_7-\tfrac43x^8-\tfrac{29}{6}y^8-\tfrac{149}{9}x^9-\tfrac{31}{3}y^9
\]
の $\kk[[x,y]]/(F_1(0),F_2(0))^2$ における class である。

\subsection{全 potential 障害加群の定義}

$P=S/I^2$ と置くと、相対偏微分は $B$-線形写像
\[
 D:P\longrightarrow T^2,\qquad [h]\longmapsto(h_x,h_y)\bmod I
\]
を定める。$d(I^2)\subset I\,dx+I\,dy$ により well-defined であり、
$\coker D=Q$ である。
中心の十九次の単位 minor を消去し、残る列のうち定数と $W_9$ を先に取ると
\begin{equation}\label{eq:Schur-obstruction}
 D\sim I_{19}\oplus[\,0\ b\ C_M\,],\qquad
 C_M:B^{12}\longrightarrow B^3.
\end{equation}
以下では
\[
 M=\coker C_M,\qquad \beta=[b]\in M,
 \qquad Q\simeq M/B\beta
\]
と書く。点族の普遍 thin module $\mathcal E$ とは異なる加群である。

\subsection{有限計算の入力とその役割}

\begin{proposition}[固定する有限入力]\label{prop:inputs}
上の族について以下が成立する。
\begin{enumerate}[label=(\alph*)]
\item $T$ は $B$ 上有限自由階数十一、$P=S/I^2$ は有限自由階数三十三である。
全空間 $T$ は正則局所整域で、$T/B$ は generically etale である。
\item 六関数 $(F_1,F_2,F_{1x},F_{1y},F_{2x},F_{2y})$ の原点の一次項は
式\eqref{eq:linear-six}で与えられる。
\item $A_0=\kk[[t_i:i\notin\{5,15,16\}]]$、$\mathfrak a=\mathfrak m_{A_0}B$ と置く。
$M,Q$ は $A_0$ 上有限、最小五元生成である。
$A_1=A_0/\mathfrak m_{A_0}^2$ に対して
\begin{equation}\label{eq:first-neighbourhood}
 M/\mathfrak a^2M\simeq Q/\mathfrak a^2Q
 \simeq A_1^4\oplus A_1/(\ell_1,\ell_2,\ell_3),
\end{equation}
\[
 \ell_1=t_8+9t_9-4t_{10}+5t_{14},\qquad
 \ell_2=t_{11}-t_{14},\qquad\ell_3=t_{13}-t_{14}.
\]
さらに $\beta\in\mathfrak n^6M\subset\mathfrak a^2M$ である。
\item 定数 class を除いた、全パラメータ四次まで延長可能な中心 critical class の空間は
$\kk[W_9]$ の一方向に限られる。この class は $B/\mathfrak n^6$ 上へ延長できる。
\end{enumerate}
\end{proposition}

\begin{proof}[入力の再現と全次数への移行]
(a) の有限平坦性は点族の determinantal 構成と長さ十一の局所完全交叉から従う。
$I/I^2\simeq T^2$ より $P$ の階数は $11+22=33$ である。
式\eqref{eq:linear-six}では $F_1,F_2$ の基底方向微分が独立であり、
形式陰関数定理により $T$ は十六次元正則局所環となる。
$B$ 上有限自由な整域であるため函数体拡大は有限で、標数零なので分離的である。
従って generically etale となる。

(b) と後に使う閉ファイバーの三次係数は
\nolinkurl{code/corank2_support_v22.py} が直接検証する。

(c) の有限性は三方向ファイバーの長さ五と
$\mathfrak n^3M\subset\mathfrak aM$ を用いた完備 Nakayama による。
後者を二回用いると $\mathfrak n^6M\subset\mathfrak a^2M$ となる。
従って十三変数側の次数高々一、全次数高々五の五百十一個の単項式で
式\eqref{eq:first-neighbourhood}を厳密に計算できる。
実際の $65\times257$ 関係行列の像は
$\ell_1h_4,\ell_2h_4,\ell_3h_4$ の三列であり、余核の長さは
$5\cdot14-3=67$ である。
\nolinkurl{archive/v21/code/relative_first_neighbourhood_v21.py} は
元の関係式、相対微分、Schur 消去への全係数再代入でこれを検証する。
(d) は中心 class の全四次条件の核と、五次までの正の相対解の直接計算である。
入口は \nolinkurl{relative_critical_quartic.py} と
\nolinkurl{relative_critical_quintic.py} であり、いずれも
\nolinkurl{archive/v21/code/} に収録する。
後者の正の解が $\beta\in\mathfrak n^6M$ を与え、
$M/\mathfrak a^2M\simeq Q/\mathfrak a^2Q$ が従う。
\end{proof}

後の議論に必要な一次項を明示する。
\[
 \lambda=t_5+t_8+9t_9-2t_{10}+4t_{11}+t_{13}+t_{14}-11t_{15},
\]
\[
 \mu=-3t_5+t_8+9t_9-10t_{10}+8t_{11}-7t_{13}
       +t_{14}-7t_{15}-17t_{16}.
\]
原点の極大イデアルを $\mathfrak m_S=(t_1,\ldots,t_{16},x,y)$ とすると
\begin{equation}\label{eq:linear-six}
\begin{array}{c|c}
\text{関数}&\mathfrak m_S/\mathfrak m_S^2\text{ における像}\\\hline
F_1&-t_{16}\\
F_2&(\lambda-17t_{16})/8\\
F_{1x}&t_{11}\\
F_{1y}&-\lambda/2-t_{15}+17t_{16}/2\\
F_{2x}&\mu/20\\
F_{2y}&t_{11}
\end{array}
\end{equation}

\section{相対 Serre 双対性による第六式の全次数消滅}\label{sec:serre-vanishing}

この節では、有限 jet の消滅とは独立に
\[
 F_{2y}\in(F_1,F_2,F_{1x},F_{1y},F_{2x})B[[x,y]]
\]
を証明する。必要なのは自己 Ext 複体の鎖自己双対性であり、
巡回 $A_\infty$ 構造や相対 potential の存在は仮定しない。

\subsection{全族における三次 Ext の自由性}

$\mathcal A_B$ を第\ref{sec:input}節の有限自由な関係代数族とし、
$\mathcal A_T=\mathcal A_B\otimes_BT$ と置く。
点族の普遍 thin module を $\mathcal E$ と書く。
各頂点で $\mathcal E_i=T$ であり、矢の作用は点族の座標である。
特に $x_3=u_{33}=z_3=1$ と正規化されている。

関係スキームの $A_*$ における完成は syzygy-dual incidence の
共通スカラーを正規化した完成と同型である。
従って、関係の基底を変えることにより、全 $B$ 上で
\begin{equation}\label{eq:relative-dual-incidence}
 L(x)=\sum_{i=1}^3x_iL_i,\quad
 Q(z)=\sum_{j=1}^3z_jQ_j,\qquad
 L_iQ_j^t=Q_jL_i^t
\end{equation}
と書ける。$L_i,Q_j$ は $6\times7$ 行列である。
これは有限次数の表示ではなく、incidence の多項式恒等式の引戻しである。

この表示の存在を確認する有限データを記す。
左 Grassmann chart の次元は七十七であり、
\[
 Q\longmapsto (L_iQ^t-QL_i^t)_{i=1}^3
\]
は $A_*$ で階数三十九、核の次元三、余核の次元六を持つ。
左関係を変分した法線写像は階数十八である。
従って三つの $Q_j$ を加えた二百三変数、百三十五方程式の incidence は
滑らかな六十八次元となる。
関係スキームへの接写像の階数は六十七で、核は $Q_j$ の共通スカラー一方向である。
関係スキーム自体の接次元も $154-87=67$ なので、
共通スカラーを切った射は両側の滑らかな六十七次元形式近傍を同型にする。
指定した十六変数 slice はこの近傍に含まれる。
これらは \nolinkurl{archive/v21/code/audit_symmetric_incidence.py} が
十四関係から有理数で再構成する。

Koszul 双加群分解を $T$ へ基底変更し、$\mathcal E$ をテンソルしてから
$\Hom_{\mathcal A_T}(-,\mathcal E)$ を適用すると、
\begin{equation}\label{eq:universal-self-ext-complex}
 \mathcal C=\operatorname{RHom}_{\mathcal A_T}(\mathcal E,\mathcal E)
 \simeq
 [\,T^4\xrightarrow{d^0}T^{12}\xrightarrow{d^1}T^{14}
       \xrightarrow{d^2}T^6\,]
\end{equation}
を得る。次数は零から三である。
この分解の完全性は特殊ファイバーの Koszul 分解と、その有限自由な相対持上げによる。
最後の六つの syzygy は式\eqref{eq:relative-dual-incidence}の六列であり、
対応する基底では
\begin{equation}\label{eq:relative-top-differential}
 d^2=[\,Q(z)\mid-L(x)\,].
\end{equation}

\begin{lemma}[Serre 固定性を用いない上端の計算]\label{lem:ext3-free}
全点族 $T$ 上で
\[
 d^2\sim\operatorname{diag}(I_5,0),\qquad
 \Ext^3_{\mathcal A_T}(\mathcal E,\mathcal E)\simeq T.
\]
この同型の自由性と三次 Ext の形成は任意の基底変更と可換する。
\end{lemma}

\begin{proof}
普遍点の左右の関係は $L(x)^tu=Q(z)^tu=0$ である。
従って $u^td^2=0$ となる。$u_{33}=1$ なので $u$ は primitive であり、
可逆な行操作によって $d^2$ の最後の行を厳密に零にできる。
一方、閉点において $d^2$ は五次の単位小行列を持つ。
例えば後述する元の関係座標では、行 $0,1,2,3,4$ と、
右関係の列 $1,3,4,5,6$ の小行列式は一である。
この小行列式は全局所環 $T$ でも単元である。
従って残る五行は列操作で $[I_5\ 0]$ となり、表示された標準形を得る。
これは複体の上端での単位消去なので基底変更後にも同じ標準形を保つ。
この議論には Serre 同型も potential も用いていない。
\end{proof}

\subsection{相対 Serre 固定性の持上げ補題}

有限自由かつ相対的に homologically smooth な代数に対する Serre 双対性は、
係数環に関する双対双加群から得られる。
体上での標準的な構成については \cite{Shklyarov} を参照されたい。
本稿で使う相対形は、有限射影分解の各項で tensor--Hom 随伴を取ることにより
同じように証明できる。
具体的に $T$-perfect な $\mathcal A_T$-perfect 対象 $E,F$ について
\begin{equation}\label{eq:relative-serre-duality}
 \operatorname{RHom}_{\mathcal A_T}(F,\mathsf S_T E)
 \simeq
 \operatorname{RHom}_T
  \bigl(\operatorname{RHom}_{\mathcal A_T}(E,F),T\bigr)
\end{equation}
であり、この同型は係数環の導来基底変更と可換する。

\begin{lemma}[上端の自由性からの Serre 固定性]\label{lem:relative-serre-fixed}
$(R,\mathfrak m,k)$ を Noether 局所環、$\mathcal A$ を上の相対 Serre 双対性を持つ
有限自由 $R$-代数とする。$E$ は各頂点で $R$-自由な有限射影次元の加群とし、
閉ファイバー $E_0$ について
\[
 \mathsf S_k(E_0)\simeq E_0[d],\qquad \End(E_0)=k
\]
と仮定する。さらに自己 Ext を計算する有限自由複体が次数 $0,\ldots,d$ に集中し、
その上端が単位消去によって自由階数一の零微分の項へ縮約されると仮定する。
このとき $\mathsf S_R(E)\simeq E[d]$ である。
\end{lemma}

\begin{proof}
$N=\mathsf S_R(E)[-d]$ と置く。
これは各頂点で $R$-perfect であり、閉ファイバーは $E_0$ の次数零の加群である。
局所環上で有限自由複体の単位成分を消去して最小複体を取ると、
剰余体上の微分は零となる。
閉ファイバーが次数零にのみ存在することから、最小複体は次数零の自由加群だけを持つ。
従って $N$ は実際の加群であり、各頂点で $E$ と同じ自由階数を持つ。

$C=\operatorname{RHom}_{\mathcal A}(E,E)$ と書く。
式\eqref{eq:relative-serre-duality}により
\[
 \Hom_{\mathcal A}(E,N)
 =H^0(C^\vee[-d])
 =\Hom_R(H^d(C),R)\simeq R.
\]
中央の等号は $C$ が次数 $0,\ldots,d$ に集中することによる。
上端の単位消去の仮定から、この $H^0$ の同定は剰余体への基底変更とも可換する。
従って左辺の生成元 $\theta:E\to N$ の閉ファイバーは
非零の $E_0\to E_0$ であり、$\End(E_0)=k$ により同型である。
各頂点で同じ階数の有限自由加群間の写像なので、Nakayama の補題によって
$\theta$ 自身も同型となる。
\end{proof}

この補題を $\mathcal E$ に適用するための閉点の Serre 条件は、
次の有限計算で検証される。残留 thin module を $E_0$ とすると、
$x=z=(0,0,1)$、$u=(0,0,0,0,0,1)$ である。
頂点 $j$ の射影加群を $P_j$ とし、$\Hom(P_j,N)=N_j$ の規約を取る。
十四関係から構成した $\operatorname{RHom}(E_0,P_j)$ のデータは
\[
\begin{array}{c|rrrr|rrr}
j&\dim C^0&\dim C^1&\dim C^2&\dim C^3
 &\rk d^0&\rk d^1&\rk d^2\\\hline
0&33&129&203&108&33&96&107\\
1&18&72&119&66&18&54&65\\
2&4&15&28&18&4&11&17\\
3&1&3&7&6&1&2&5
\end{array}
\]
である。各行のコホモロジーは $(0,0,0,1)$ となる。
さらに射影加群間の十二の prefix 写像が三次 Ext に誘導する作用を双対化すると、
$\mathsf S(E_0)[-3]$ の全矢は $E_0$ の全矢と一致する。
頂点の基底変換は四つとも一に取れる。
\nolinkurl{archive/v21/code/point_serre.py --fat} はこの階数、微分の合成、
十二の矢を全て厳密に再計算する。
従って $\mathsf S(E_0)\simeq E_0[3]$ は次元比較だけの主張ではない。
また $\End(E_0)=\kk$、自己 Ext の次元は $(1,2,2,1)$ である。

\begin{corollary}\label{cor:universal-serre-fixed}
普遍 thin module は全 $T$ 上で
\[
 \mathsf S_T(\mathcal E)\simeq\mathcal E[3]
\]
を満たす。従って $\mathcal C\simeq\mathcal C^\vee[-3]$ である。
\end{corollary}
\begin{proof}
補題\ref{lem:ext3-free}と補題\ref{lem:relative-serre-fixed}を適用する。
最後の同型は式\eqref{eq:relative-serre-duality}にこの Serre 同型を代入したものである。
\end{proof}

\subsection{自己 Ext 複体と実際の二変数 Jacobian}

\begin{proposition}[単位消去による最小複体]\label{prop:minimal-self-ext}
式\eqref{eq:universal-self-ext-complex}は $T$ 上で
\begin{equation}\label{eq:minimal-universal-ext}
 [\,T\xrightarrow{0}T^2\xrightarrow{J_F}T^2\xrightarrow{0}T\,],
 \qquad J_F=\begin{pmatrix}F_{1x}&F_{1y}\\F_{2x}&F_{2y}\end{pmatrix},
\end{equation}
とホモトピー同値である。特にある $U,V\in\GL_2(T)$ に対して
\begin{equation}\label{eq:chain-self-duality}
 UJ_F=J_F^tV
\end{equation}
が全 $T$ 上で成立する。
\end{proposition}

\begin{proof}
最初の微分は頂点基底変換の無限小作用である。
$x_3,u_{33},z_3$ の三つの単元を用いて、次数零と一の間の三つの可縮対を除く。
次数零には共通スカラー $T$ だけが残り、その微分は零である。
次数一の残る九座標は
\[
 x,y,u_{11},u_{12},u_{13},u_{22},u_{23},z_1,z_2
\]
であり、$d^1$ は十四関係のこれらに関する Jacobian となる。

次数二と三の五つの可縮対を、二変数表示と両立する形で取り除く。
元の右関係行列を $R$ とすると、全族上で $Q=RA(t)$ を与える可逆行列を取れる。
その閉点での値は
\[
A(0)=\begin{pmatrix}
0&-1&0&0&0&0&0\\
-1&1&0&1&0&0&0\\
0&0&0&-1&0&0&0\\
0&-1&-1&1&0&-1&0\\
0&0&0&0&0&-1&0\\
0&0&0&-1&-1&0&0\\
0&0&0&0&0&0&1
\end{pmatrix},\qquad\det A(0)=-1.
\]
従って元の左右の関係座標では $d^2=[Q,-LA(t)^t]$ である。
行 $0,1,2,3,4$、右関係の列 $1,3,4,5,6$ の小行列式は閉点で一となる。
補題\ref{lem:ext3-free}の primitive 左核により $\operatorname{im}d^2$ の階数は五である。
従って $\ker d^2$ は、全七左関係と右関係 $0,2$ の九座標への射影で $T^9$ と同型になる。
この射影を使って五つの可縮対を除去してよい。

この九関係のうち、左関係 $0,2,4,5,6$ と右関係 $0,2$ の七本を取り、
最後の七変数 $u_{11},\ldots,u_{23},z_1,z_2$ で微分した小行列を考える。
左の $5\times5$ block、右の $2\times2$ block の行列式は閉点でともに一であり、
全 $7\times7$ block は下三角 block なのでその行列式も一である。
従って七つの補助変数を形式陰関数定理で消去できる。
Jacobian の Schur complement は、消去後の残り二関係を $x,y$ で微分した行列に一致する。
この操作は次数一と二の七つの可縮対の除去にほかならない。

残る二関係は左の行 $3,1$ であり、選んだ $F_1,F_2$ はこれらを
共通単元 $-\det C$ 倍したものである。ここで $C$ は五つの $u$ を消去する小行列である。
$T$ 上では関係そのものが零なので、共通単元の微分による項は消える。
従って可逆な行基底変更をして、中間微分を正確に $J_F$ と取れる。
これで式\eqref{eq:minimal-universal-ext}が得られる。
相対自己 Ext と接空間の次元の一致だけからこの表示を推定したのではなく、
実際の射影分解の三、五、七個の単位消去で得たことに注意する。

系\ref{cor:universal-serre-fixed}の鎖同型をこの縮約へ移す。
式\eqref{eq:minimal-universal-ext}の微分は局所環 $T$ の極大イデアルに属するので、
この複体も双対複体も最小である。
有界有限自由最小複体のホモトピー同値は、剰余体への基底変更で各次数の同型を与え、
従って各次数で可逆行列により表される。
次数一、二の鎖写像条件が式\eqref{eq:chain-self-duality}である。
双対の符号規約によって負号が現れる場合には $V$ に吸収する。
\end{proof}

\subsection{二次初項が強制する off-diagonal 単元}

\begin{lemma}\label{lem:offdiagonal-unit}
式\eqref{eq:chain-self-duality}を満たす可逆行列 $U,V$ の定数項は
\[
 U(0)=V(0)=c\begin{pmatrix}0&1\\1&0\end{pmatrix},\qquad c\ne0.
\]
従って $U_{12}$ は $T$ の単元である。
\end{lemma}

\begin{proof}
基底パラメータ $t$ を零にした後、$x,y$ の二次項を比較する。
中心の関係は三次以上なので、$x^2,xy,y^2$ は商環の二次部分でも独立である。
また
\[
 J_2=\begin{pmatrix}2xy&x^2\\y^2&2xy\end{pmatrix}.
\]
従って $U(0)J_2=J_2^tV(0)$ である。
$xy$ の係数は $U(0)=V(0)$ を与える。
$x^2$ の係数は $U_{11}(0)=V_{11}(0)=0$、$U_{21}(0)=V_{12}(0)$ を、
$y^2$ の係数は $U_{22}(0)=V_{22}(0)=0$、$U_{12}(0)=V_{21}(0)$ を与える。
従って表示された形に限られ、可逆性から $c\ne0$ となる。
\end{proof}

\begin{theorem}[第六 Jacobian 成分の全次数所属]\label{thm:psi-vanishing}
実際の十六変数族に対して
\[
 \psi=[F_{2y}]=0\quad\text{in}\quad
 \Lambda=B[[x,y]]/(F_1,F_2,F_{1x},F_{1y},F_{2x}).
\]
従って corank 二の環は五つの式だけで定義される。
\end{theorem}

\begin{proof}
$T$ から $\Lambda$ へ移すと
$J_F=\left(\begin{smallmatrix}0&0\\0&\psi\end{smallmatrix}\right)$ となる。
式\eqref{eq:chain-self-duality}の $(1,2)$ 成分は $U_{12}\psi=0$ である。
補題\ref{lem:offdiagonal-unit}により $U_{12}$ は単元なので $\psi=0$ を得る。
\end{proof}

\subsection{全形式の相対 almost closed 一形式}

相対一形式 $\alpha=g_1\,dx+g_2\,dy\in\Omega^1_{S/B}$ が
\emph{almost closed} であるとは、
\[
 d_{S/B}\alpha\in(g_1,g_2)\Omega^2_{S/B}
\]
を満たすこと、すなわちその外微分が自身の零スキーム上で消えることとする。
ここでも微分は $x,y$ に関する連続相対微分である。

\begin{proposition}[相対 almost closed 表示]\label{prop:relative-almost-closed}
ある $\widetilde K\in\GL_2(S)$ を用いて
\[
 \binom{g_1}{g_2}=\widetilde K\binom{F_1}{F_2},\qquad
 \alpha=g_1\,dx+g_2\,dy
\]
と置くと、$\alpha$ は全形式の相対 almost closed 一形式となる。
その係数のイデアルは $I=(F_1,F_2)$ に一致し、零スキームは元の点族 $\Spec T$ である。
さらに $\widetilde K$ と $\alpha$ の係数を代数的冪級数として選べる。
従って複素解析的な原点近傍で収束する相対 almost closed 表示も存在する。
\end{proposition}

\begin{proof}
式\eqref{eq:chain-self-duality}とその転置を合わせ、
\[
 K=\frac{U+V^t}{2}\in\operatorname{Mat}_2(T)
\]
と置く。標数零なのでこの平均を取ることができ、
\[
 KJ_F=\frac{UJ_F+V^tJ_F}{2}
      =\frac{J_F^tV+J_F^tU^t}{2}
      =J_F^tK^t=(KJ_F)^t
\]
となる。補題\ref{lem:offdiagonal-unit}により
$K(0)=c\left(\begin{smallmatrix}0&1\\1&0\end{smallmatrix}\right)$、$c\ne0$ なので
$K$ は可逆である。
$S\to T$ に沿って各成分を持ち上げれば、行列式の定数項が非零であるため
$\widetilde K\in\GL_2(S)$ を取れる。

$g=\widetilde K F$ の Jacobian を $I$ を法として計算すると、
$\widetilde K$ の微分に $F$ を掛けた項は消えるため
$J_g\bmod I=KJ_F$ となる。
これは対称行列なので $g_{2x}-g_{1y}\in I$、従って
$d_{S/B}\alpha\in I\Omega^2_{S/B}$ である。
また $\widetilde K$ の可逆性から $(g_1,g_2)=I$ が成立する。

代数性も有限線形系から得られる。
第\ref{sec:input}節の relation slice の係数は代数的冪級数なので、
\[
 S_{\mathrm{alg}}=\kk\langle t_1,\ldots,t_{16},x,y\rangle
 =\bigl(\kk[t_1,\ldots,t_{16},x,y]_{(t,x,y)}\bigr)^h,
 \qquad R=S_{\mathrm{alg}}/(F_1,F_2)
\]
を取れる。$R$ は Noether 局所環で、その完備化は $T$ であり、
$R\to T$ は忠実平坦である \cite[Tags 00MC,0A1W]{Stacks}。
上で得た $K$ を $c$ で割り、原点値を
$\Sigma=\left(\begin{smallmatrix}0&1\\1&0\end{smallmatrix}\right)$ に規格化する。
$R$ の極大イデアルの有限生成元を $m_1,\ldots,m_e$ とし、
\[
 K=\Sigma+\sum_{i=1}^e m_iH_i
\]
を条件 $KJ_F-J_F^tK^t=0$ に代入する。
これは未知行列 $H_i$ の成分に関する $R$ 上の有限個の非斉次線形方程式である。
完備化 $T$ 上では解を持つため、その線形写像の余核における右辺の class は零となる。
忠実平坦性により同じ class は $R$ 上でも零であり、$H_i$ を $R$ 上で選べる。
この解は原点値 $\Sigma$ を保つので、得られた $K\in\GL_2(R)$ を
$\widetilde K\in\GL_2(S_{\mathrm{alg}})$ へ持ち上げられる。
前半と同じ微分計算により代数的な almost closed 表示を得る。
複素数上の代数的冪級数は原点近傍で収束するため、最後の主張も従う。
ここでは既存の任意の形式表示の収束性を主張せず、収束する表示を選んでいる。
\end{proof}

almost closed 条件は $d_{S/B}\alpha$ が零スキーム上で消える条件であり、
$S$ 上の閉性 $d_{S/B}\alpha=0$ はこの証明からは従わない。
一般に almost closed 一形式の零スキームが局所的な potential の critical locus に
なるとは限らず、Pandharipande--Thomas はその反例を構成している
\cite{PandharipandeThomasAlmostClosed}。
本例における全相対 potential の存在、すなわち $\beta=0$ は引き続き未決定であり、
その正確な条件を第\ref{sec:obstruction}節の一つの元のイデアル所属へ帰着する。

\section{正則な corank 二の環と四枚被覆}\label{sec:fourfold}

\subsection{五つの式の消去}

\begin{definition}
次の完備局所環とその元を取る。
\begin{align*}
\Lambda&=B[[x,y]]/(F_1,F_2,F_{1x},F_{1y},F_{2x}),\\
\psi&=[F_{2y}]\in\Lambda.
\end{align*}
定理\ref{thm:psi-vanishing}により $\psi=0$ である。
従って $\Lambda=T/\mathfrak j$ は内在的な corank 二の環そのものである。
\end{definition}

\begin{theorem}[四枚被覆]\label{thm:fourfold}
$\Lambda$ は十三次元正則局所整域であり、$A_0$ 上有限自由階数四である。
さらに
\begin{equation}\label{eq:lambda-fiber}
\Lambda/\mathfrak m_{A_0}\Lambda
\simeq\kk[[x,y]]/(xy,y^2-2x^2),
\end{equation}
であり、このファイバーで
\[
t_5=11x^2,\qquad t_{15}=x^2,\qquad t_{16}=0.
\]
$1,x,y,x^2$ は $\Lambda$ の $A_0$-基底に取れる。
従って
\[
\Omega^2_{T/B}\simeq\Lambda\,dx\wedge dy.
\]
\end{theorem}

\begin{proof}
最初の五式の $(t_5,t_8,t_{11},t_{15},t_{16})$ に関する一次微分行列は
\[
\begin{pmatrix}
0&0&0&0&-1\\
1/8&1/8&1/2&-11/8&-17/8\\
0&0&1&0&0\\
-1/2&-1/2&-2&9/2&17/2\\
-3/20&1/20&2/5&-7/20&-17/20
\end{pmatrix}.
\]
行列式は $1/40$ である。従って形式陰関数定理により、
この五変数は他の十一個の $t$ と $x,y$ の形式級数として一意に解ける。
$\Lambda$ は十三変数の形式冪級数環であり、特に正則局所整域である。

$\mathfrak m_{A_0}$ で割り、$a=t_5,b=t_{15},c=t_{16}$ と書く。
五式が定める $\kk[[a,b,c,x,y]]$ の商について、
全次数 $N$ までの Macaulay 行列を作る。
実際の係数から得られる値は
\[
\begin{array}{c|c|c|c}
N&\text{行列の大きさ}&\text{階数}&\text{余核次元}\\\hline
1&6\times4&3&3\\
2&21\times25&17&4\\
3&56\times90&52&4
\end{array}
\]
である。最後の二つの長さの一致は、商環の極大イデアル $\mathfrak m$ に対して
$\mathfrak m^3/\mathfrak m^4=0$ を意味する。
Nakayama の補題により $\mathfrak m^3=0$ であり、全商環の長さは四である。
三次余核の基底は $1,x,y,x^2$、全座標の作用は
\[
xy=0,\quad y^2=2x^2,\quad
a=11x^2,\quad b=x^2,\quad c=0
\]
となる。これらの式を全入力へ再代入し、四基底の独立性も確認する。
これで式\eqref{eq:lambda-fiber}が全次数で従う。

この有限長ファイバーにより
$\mathfrak m_\Lambda^3\subset\mathfrak m_{A_0}\Lambda$ である。
両位相は同値なので、完備版 Nakayama の議論によって
$1,x,y,x^2$ が $\Lambda$ を $A_0$ 上生成する。
$A_0$ と $\Lambda$ はともに十三次元正則局所環で、閉ファイバーは零次元である。
したがって miracle flatness \cite[Tag 00R4]{Stacks} により $\Lambda/A_0$ は平坦である。
有限平坦加群は自由であり、閉ファイバーの長さから階数は四である。
四元生成の全射 $A_0^4\to\Lambda$ は閉ファイバーで同型なので同型となる。
最後の主張は命題\ref{prop:omega2}と定理\ref{thm:psi-vanishing}による。
\end{proof}

\section{四枚被覆の全形式標準形と $S_4$}\label{sec:four-versal-v23}

この節では、選んだ十三変数の基底 $A_0$ 上の有限代数 $\Lambda$ を調べる。
四枚被覆としての構造は、四変数の最小versal変形と九変数の自明な方向に分離できる。
ここで変形とは、特殊ファイバーの同一視を持つ有限平坦可換代数の変形を指す。
各変形に原点の切断を要求しない。また、以下の標準形を与える座標と分類写像は一意ではない。

\subsection{長さ四の完全交叉の最小versal変形}

\begin{lemma}\label{lem:four-miniversal-v23}
標数零の体 $k$ 上で
\[
 D_0=k[[X,Y]]/(XY,Y^2-2X^2)
\]
と置く。$D_0$ の可換代数変形の接空間は四次元である。
次の族は最小versal変形であり、その全空間は四次元正則局所環である。
\begin{align}
 R&=k[[u,v,w,z]],\\
 \mathscr D&=R[[X,Y]]/(XY-u,\ Y^2-2X^2-v-wX-zY).
 \label{eq:four-miniversal-v23}
\end{align}
$\mathscr D$ は $R$ 上有限自由階数四である。
\end{lemma}

\begin{proof}
完全交叉の余接複体を使うと、接空間は
\[
 T^1(D_0)=\coker\left[
 D_0^2\xrightarrow{
 \left(\begin{smallmatrix}Y&X\\-4X&2Y\end{smallmatrix}\right)}D_0^2
 \right].
\]
この余核では $(X,0)=(0,-2Y)$、$(Y,0)=(0,4X)$ であり、
二つの成分の二次項は全て消える。従って
\begin{equation}
 [(1,0)],\quad[(0,1)],\quad[(0,X)],\quad[(0,Y)]
 \label{eq:four-T1-basis-v23}
\end{equation}
が基底となる。この記述は完全交叉の余接複体の標準的な二項表示
\cite[Tag 08SH]{Stacks}による。

最小versalityを、この場合には直接確認できる。
完備Noether局所 $k$-代数 $C$ 上の有限平坦変形を考える。
$X,Y$ を持ち上げると $C[[X,Y]]$ から全射があり、
平坦性とNakayamaの補題により、その核は中心の正則列の持上げ二本で生成される。
従って変形を、持ち上げた $X,Y$ と二つの関係で表示できる。
基底の極大イデアルを $\mathfrak m_C$ とし、標準形を
$\mathfrak m_C^n$ を法として構成したとする。
次の次数の二つの誤差を $D_0^2\otimes_k
\mathfrak m_C^n/\mathfrak m_C^{n+1}$ で考え、
その $T^1$ 成分を四つの係数 $u,v,w,z$ の変更で吸収する。
残りのJacobian像は $X,Y$ の変更で消え、
二つの中心関係の倍は可逆な関係基底の変更で消える。
この操作を帰納的に続けると完備性により全形式の標準形を得る。
接写像が基底\eqref{eq:four-T1-basis-v23}を与えるので最小である。
これは一意な分類写像の主張ではない。

最後に $u,v$ を消去すると
$\mathscr D\simeq k[[X,Y,w,z]]$ である。
閉ファイバーの長さが四なので、完備Nakayamaにより有限であり、
両環の正則性と次元一致から平坦となる。
局所環上の有限平坦加群は自由で、その階数は閉ファイバーの長さ四である。
\end{proof}

\subsection{実際の族のKodaira--Spencer写像}

\begin{proposition}\label{prop:actual-four-KS-v23}
実際の $\Lambda/A_0$ のKodaira--Spencer写像は全射である。
特に $t_8,t_{11},t_{13},t_{14}$ の四方向だけで $T^1(D_0)$ を生成する。
\end{proposition}

\begin{proof}
この有限線形計算は、元の十六変数族の関係式から行う。
$a=t_5,b=t_{15},c=t_{16}$ と書き、閉ファイバーでの五式を
\[
 G=(F_1,F_2,F_{1x},F_{1y},F_{2x}),\qquad
 \mathbf q=(a,b,c,x,y)
\]
とする。$D_0$ 上では
$a=11x^2,b=x^2,c=0$ である。
五変数五式の余接表示から
\[
 T^1(D_0)=\coker\bigl(D_0^5\xrightarrow{J_{\mathbf q}G}D_0^5\bigr).
\]
基底 $(1,x,y,x^2)$ を使うと $J_{\mathbf q}G$ は二十次正方行列で、
有理数上の階数は十六である。各基底方向 $t_i$ の像は
$[\partial G/\partial t_i]$ で与えられる。
左零化空間のある基底で、上記四方向の列は
\[
 \begin{pmatrix}
 0&1&0&0\\
 1/5&1&-1/5&1/5\\
 -44/25&-51/5&-6/25&-44/25\\
 0&26/5&-4&-4
 \end{pmatrix},
 \qquad \det=-\frac85.
\]
従って接写像は全射である。

検証器\nolinkurl{four_cover_versal_v23.py}は、
実際の混合関係係数から二十次Jacobian、左零化空間、十三本の接方向を再構成し、
Fraction演算と独立な有理行列演算で階数と小行列式を検査する。
入力の高次打切りはこの計算に影響しない。
閉ファイバーで $(a,b,c)^2=0$、$(x,y)^3=0$ なので、
一回のパラメータ微分と高々二回の $x,y$ 微分の後に残り得る項は、
保存された実際の混合jetに全て含まれる。
\end{proof}

\begin{theorem}[四枚被覆の標準形]\label{thm:four-normal-form-v23}
ある形式座標と $A_0$-代数同型により
\begin{align}
 A_0&\simeq\mathbb C[[u,v,w,z,s_1,\ldots,s_9]],\\
 \Lambda&\simeq
 \frac{A_0[[X,Y]]}{(XY-u,\ Y^2-2X^2-v-wX-zY)}.
 \label{eq:actual-four-normal-form-v23}
\end{align}
従って $\Lambda/A_0$ は、式\eqref{eq:four-miniversal-v23}の族に
九変数の形式円板を加えたものである。
\end{theorem}

\begin{proof}
補題\ref{lem:four-miniversal-v23}により、分類写像 $R\to A_0$ と
$\Lambda\simeq\mathscr D\widehat\otimes_R A_0$ を取れる。
命題\ref{prop:actual-four-KS-v23}と $R$ の最小性から、
$u,v,w,z$ の像の一次項は $\mathfrak m_{A_0}/\mathfrak m_{A_0}^2$ で独立である。
これらを十三個の正則パラメータへ延ばし、形式的逆関数定理を適用すると
$A_0\simeq R[[s_1,\ldots,s_9]]$ を得る。
これは $\Lambda$ のみに関する標準形であり、
元の十一点族 $T/B$ が相対potentialを持つことは意味しない。
\end{proof}

\subsection{ramification divisorと一般Galois群}

\begin{corollary}\label{cor:four-ramification-v23}
$\Spec\Lambda\to\Spec A_0$ のramification divisorは、形式的に
三次元の普通二重点と九次元形式円板の積である。
具体的には $\Lambda\simeq\mathbb C[[X,Y,w,z,s_1,\ldots,s_9]]$ の中で
\[
 4X^2+2Y^2+wX-zY=0
 \label{eq:four-ramification-v23}
\]
と書ける。その特異部分は $X=Y=w=z=0$ で九次元である。
\end{corollary}

\begin{proof}
有限写像は $u=XY$、$v=Y^2-2X^2-wX-zY$ である。
その $X,Y$ に関するJacobian行列式は
$2Y^2-zY+4X^2+wX$ である。
ここで線形座標
\[
 \alpha=X,\qquad\beta=4X+w,\qquad
 \gamma=Y,\qquad\delta=2Y-z
\]
を取ると、式は $\alpha\beta+\gamma\delta=0$ となる。
従って三次元の普通二重点と九つの自明な方向の積である。
ここで記述しているのは被覆の全空間内のramification divisorであり、
基底内のbranch divisorそのものの標準形ではない。
\end{proof}

\begin{theorem}\label{thm:four-S4-v23}
四枚被覆 $\Lambda/A_0$ の一般ファイバーのGalois閉包の群は $S_4$ である。
\end{theorem}

\begin{proof}
まず九つの自明な変数を除いた族を考える。
$u\ne0$ では $Y=u/X$ と消去でき、$X$ の方程式は
\[
 2X^4+wX^3+vX^2+zuX-u^2=0.
 \label{eq:four-quartic-v23}
\]
完備化によってGalois群が小さくなる可能性も除くため、
順序付きの四根 $r_1,\ldots,r_4$ の代数を直接調べる。
$e_i$ を四根の基本対称式とすると、Vietaの式は
\[
 e_1=-w/2,\qquad e_2=v/2,\qquad
 2e_3+uz=0,\qquad 2e_4+u^2=0.
\]
従って $w,v$ を消去した多項式環での分解代数は
\[
 C=\frac{\mathbb C[r_1,r_2,r_3,r_4,u,z]}
 {(u^2+2e_4,\ uz+2e_3)}.
\]
最初の式は既約であり、二つの式は正則列をなす。
一方 $u=0$ のファイバーでは $e_3=e_4=0$ となる。
これは少なくとも二つの根が零という条件であり、$z$ を加えて次元三である。
$C$ は次元四の完全交叉なので、$u$ は非零因子である。
また
\[
 C[u^{-1}]
 \simeq
 \frac{\mathbb C[r_1,r_2,r_3,r_4,u,u^{-1}]}
 {(u^2+2r_1r_2r_3r_4)}
\]
は整域である。実際、四根の積は分数体内の平方ではない。
従って $C$ 自身が整域である。

$r_i,z$ の重みを一、$u$ の重みを二とすると $C$ は正次数付き整域である。
重みの位相は原点の極大イデアル進位相と同値であり、
その完備化 $\widehat C$ の随伴次数環は $C$ である。
従って $\widehat C$ も整域となる。
四次のモニック多項式の順序付き分解代数は基底上自由階数 $4!=24$ である
（根を一つずつ加え、残りのモニック因子で割る構成による）。
特殊ファイバーの全ての根は零なので、$\widehat C$ は
$\mathbb C[[u,v,w,z]]$ 上のその分解代数であり、階数二十四を持つ。
判別式が零多項式でないことは、$u=1,w=v=z=0$ で $2X^4-1$ となることから分かる。
この非零性は整域の完備化でも保たれるので、一般の四根は相異なり、
$S_4$ はそれらを忠実に置換する。
分数体の次数も二十四なのでGalois群は正確に $S_4$ である。

九つの自明な形式変数を加えた順序付き分解代数は
$\widehat C[[s_1,\ldots,s_9]]$ で、引き続き整域である。
したがって $\Lambda/A_0$ に対しても同じ結論となる。
\end{proof}

この $S_4$ は、十三変数への投影で現れる四つのcorank二点のmonodromyである。
元の残留十一点の $S_{11}$、および全二十一点の $S_{21}$ とは対象が異なる。

\section{分岐自由因子と普通二重点による正規化}\label{sec:four-branch-v24}

前節の四枚被覆について、全空間のramification divisorと基底のbranch divisorを
区別して記述する。九つの自明な変数を除き、
$R=\mathbb C[[u,v,w,z]]$、$\mathscr D\simeq\mathbb C[[X,Y,w,z]]$ とする。
この節の行列式は多項式恒等式であり、有限jetへの打切りを含まない。

\subsection{trace判別式と対数的ベクトル場}

\begin{theorem}\label{thm:four-free-discriminant-v24}
次の行列と多項式を定める。
\begin{equation}\label{eq:four-discriminant-matrix-v24}
C=\begin{pmatrix}
 uw/4&-v/2&2u&(-3uz+vw)/4\\
 3uz&4u-wz&2v&(-uw-vz+w^2z)/2\\
 (8v-w^2)/4&-3z&w&(16u+5wz)/4\\
 -4u&3w/2&z&(4v-3w^2)/4
\end{pmatrix},\qquad \delta=\det C.
\end{equation}
四枚被覆のbranch divisorは被約既約な超曲面 $\Delta=V(\delta)$ である。
$\delta$ は基底 $(1,X,Y,X^2)$ に関するtrace判別式であり、
重み $(2,2,1,1)$ について次数八の準斉次多項式である。
$-C$ の四つの列を係数とするベクトル場は
$\Der_R(-\log\Delta)$ の基底を与える。特に $\Delta$ は自由因子である。
さらに、$q=4X^2+2Y^2+wX-zY$ と置くと
\begin{align}
 \Nm_{\mathscr D/R}(q)&=-8\delta,\label{eq:four-norm-v24}\\
 \operatorname{disc}_X(2X^4+wX^3+vX^2+zuX-u^2)&=16u^2\delta.
 \label{eq:four-quartic-index-v24}
\end{align}
従って四次多項式の判別式には、四枚被覆自体の分岐に由来しない因子 $u^2$ がある。
\end{theorem}

\begin{proof}
$\mathscr D$ の有限自由基底 $(1,X,Y,X^2)$ における乗法行列は
\[
M_X=\begin{pmatrix}
0&0&u&-zu/2\\1&0&0&-v/2\\0&0&0&u/2\\0&1&0&-w/2
\end{pmatrix},\qquad
M_Y=\begin{pmatrix}
0&u&v&0\\0&0&w&u\\1&0&z&0\\0&0&2&0
\end{pmatrix}.
\]
直接乗算により $M_XM_Y=uI$、
$M_Y^2-2M_X^2=vI+wM_X+zM_Y$ が成り立つ。
$M_1=I,M_2=M_X,M_3=M_Y,M_4=M_X^2$ とすると
\[
 \det\bigl(\operatorname{Tr}(M_iM_j)\bigr)_{i,j=1}^4=\det C.
\]
同じ行列の行列式から式\eqref{eq:four-norm-v24}が得られ、
四次多項式のresultantから式\eqref{eq:four-quartic-index-v24}が得られる。
これらは有理数係数の多項式の直接計算である。
trace判別式による非\'etale locusの記述は
\cite[Tag 0BVH]{Stacks}の標準的な判定とも一致する。

行列 $C$ は相対接コホモロジーの最小表示から得られる。
\[
 T^1_{\mathscr D/R}=\coker\left[
 \mathscr D^2\xrightarrow{
 J=\left(\begin{smallmatrix}Y&X\\-4X-w&2Y-z\end{smallmatrix}\right)}
 \mathscr D^2\right]
\]
を上の基底で八次正方行列に展開する。
行 $(2,3,4,8)$、列 $(1,2,3,5)$ の小行列式は $-8$ である。
この単位ブロックを消去したSchur補行列が $C$ であり、
同じ消去の下で $u,v,w,z$ のKodaira--Spencer写像は $-I_4$ になる。

ここで被約性と既約性を確認する。
$V(q)$ は $\alpha\beta+\gamma\eta=0$ という三次元普通二重点なので整かつ正規である。
有限射によるその像は既約であり、normの等式からその台は $V(\delta)$ に等しい。
多項式族の基底点 $(u,v,w,z)=(1,-7,0,6)$ では
\[
 2X^4-7X^2+6X-1=(X-1)^2(2X^2+4X-1).
\]
唯一の重根は $X=1$ であり、対応する点は $(X,Y)=(1,1)$ である。
ここで $\ker J$ は $(1,-1)$ で生成され、
$dq(1,-1)=10$ だから単純なfoldである。
従って像の一般点における判別式の重複度は一である。
形式完備化の後も $V(q)$ は整であり、有限射の像の台は既約である。
また準斉次多項式 $\delta$ の被約性は完備化で失われない。
実際、形式級数の非単元の二乗が $\delta$ を割れば、最低重み項の比較により
多項式 $\delta$ にも非定数の平方因子が生じる。これで形式的にも被約既約性が従う。

最後に、直接微分によって
\begin{equation}\label{eq:four-saito-identity-v24}
 (\partial_u\delta,\partial_v\delta,\partial_w\delta,\partial_z\delta)(-C)
 =\delta\,(w/2,\,0,\,-8,\,3z/2)
\end{equation}
を得る。$\det(-C)=\delta$ と被約性にSaitoの判定を適用すればよい
\cite[Propositions 1.3, 1.5]{BEvB}。
\end{proof}

自由性そのものは、最小versal有限平坦写像に関する一般定理の適用でもある
\cite[Theorems 3.1, 3.4]{BEvB}。
ここでは本研究のcorank二の被覆について、trace、different、最小表示、
対数的ベクトル場を同じ四次行列で具体化した。
式\eqref{eq:four-quartic-index-v24}の因子 $u^2$ は、$u=0$ で
$X$ 一つでは有限代数の基底を与えられなくなることに対応する。
これを除かない四次判別式をbranch divisorの被約方程式とすることはできない。

\subsection{分岐因子の正規化}

\begin{theorem}\label{thm:four-branch-normalization-v24}
有限射
\[
 V(q)\longrightarrow\Delta
\]
は正規化写像である。従ってbranch divisorの正規化は三次元普通二重点である。
さらに $N=T^1_{\mathscr D/R}$ を有限 $R$-加群と見れば
\begin{equation}\label{eq:four-fast-normalization-v24}
 \End_R(N)\simeq\mathscr D/(q)
 \simeq\mathbb C[[\alpha,\beta,\gamma,\eta]]/(\alpha\beta+\gamma\eta).
\end{equation}
実際の十三変数の基底では、これらの環と写像に九変数の形式円板を加える。
\end{theorem}

\begin{proof}
被約既約な素因子 $(\delta)$ で $R$ を局所化すると離散付値環となる。
有限自由加群 $\mathscr D$ 上の $q$ 倍写像の行列式は $-8\delta$ なので、
その余核 $(\mathscr D/(q))_{(\delta)}$ の長さは一である。
この余核は $\Frac(R/(\delta))$ 上有限な整域、従って体であり、
その長さは $\Frac(\mathscr D/(q))$ の拡大次数に等しい。
従って有限射 $V(q)\to\Delta$ は形式環上でも一般次数一である。
$V(q)$ は正規であるから、この有限双有理射は正規化写像である。

$\det J=q$ と余因子行列の恒等式から $qN=0$ である。
さらに余因子行列は
\[
 N\longrightarrow(\mathscr D/(q))^2,\qquad
 [a]\longmapsto\operatorname{adj}(J)a\bmod q
\]
という単射を与える。実際、$\operatorname{adj}(J)a=qb$ なら
$q a=qJb$ となり、$\mathscr D$ が整域なので $a=Jb$ である。
$J$ は $\Frac(\mathscr D/(q))$ 上で階数一を持つため、
$N$ は $\mathscr D/(q)$ 上の階数一のtorsion-free加群となる。
従って $R/(\delta)$ 上でも階数一かつtorsion-freeであり、
その自己準同型環は共通の関数体の部分環である。
乗法によって $\mathscr D/(q)\subset\End_R(N)$ があり、
他方 $\End_R(N)$ は $R/(\delta)$ 上有限なので、その全ての元は整である。
既に $\mathscr D/(q)$ が整閉包と分かったので、逆包含も従う。
\end{proof}

式\eqref{eq:four-fast-normalization-v24}は、この例における
接コホモロジー加群からの正規化の具体的な実現である。
この加群 $N$ は、次節以降の元の十一点族のpotential障害加群 $M$ とは異なる。
検証器\nolinkurl{four_cover_discriminant_v24.py}は、
乗法関係、trace判別式、Schur消去、式\eqref{eq:four-saito-identity-v24}、
normと四次判別式の係数、foldの横断微分を有理数上で再計算する。

\section{四枚被覆の非対称性と線形pencilの critical slice}
\label{sec:four-noncritical-v24}

この節では、四枚被覆そのものを相対 critical locus として表示できるかを調べる。
対象は四変数の最小versal変形
\begin{equation}\label{eq:four-noncritical-input-v24}
 R=\mathbb C[[u,v,w,z]],\qquad
 D=\frac{R[[X,Y]]}{(f_1,f_2)},\qquad
 f_1=XY-u,\quad f_2=Y^2-2X^2-v-wX-zY
\end{equation}
である。$u,v$ を消去すると $D\simeq\mathbb C[[X,Y,w,z]]$ となる。
$D/R$ は有限自由階数四であり、相対Jacobianは
\begin{equation}\label{eq:four-jacobian-v24}
 J=\begin{pmatrix}Y&X\\-4X-w&2Y-z\end{pmatrix},\qquad
 \det J=4X^2+2Y^2+wX-zY
\end{equation}
である。特に $D$ は整域であり、$D/R$ はgenerically etaleである。

結論は、四変数族全体では相対対称obstruction theoryさえ存在しない一方、
$(w,z)$ 平面内のほとんどすべての一次元方向では三変数の多項式potentialが存在する、
というものである。例外は明示的な二方向だけである。
以下では完備相対余接複体を $\mathbb L_{D/R}$ と書く。
本例では有限代数の余接複体と一致し、完全交叉の二項表示で計算できる
\cite[Tag 08SH]{Stacks}。

\subsection{相対対称obstruction theoryに対する必要条件}

相対perfect obstruction theoryとは、perfect amplitudeが $[-1,0]$ に含まれる
複体 $E$ と射 $\phi:E\to\mathbb L_{D/R}$ で、
$h^0(\phi)$ が同型、$h^{-1}(\phi)$ が全射となるものをいう。
これが対称であるとは、同型
\[
 \theta:E\xrightarrow{\sim}E^\vee[1],\qquad
 \theta^\vee[1]=\theta
\]
を持つことをいう。これは通常の対称perfect obstruction theoryの定義の相対版である
\cite[Definition 1.10]{BehrendFantechiSymmetric}。

\begin{lemma}\label{lem:symmetric-pot-is-cotangent-v24}
$R$ を正則完備局所環、$D$ を整域である有限平坦局所 $R$-代数とし、
$D/R$ は局所完全交叉かつgenerically etaleであるとする。
このとき相対対称perfect obstruction theory
$E\to\mathbb L_{D/R}$ が存在すれば、この射は擬同型である。
従って必要条件として
\[
 \mathbb L_{D/R}\simeq\mathbb L_{D/R}^\vee[1]
\]
を満たす。
\end{lemma}

\begin{proof}
$D$ は局所環なので、$E$ を有限自由加群の二項複体
$[F^{-1}\to F^0]$ で表せる。
対称性から $\operatorname{rk}F^{-1}=\operatorname{rk}F^0$ である。
$K=\operatorname{Frac}(D)$ と置くと、generically etaleという仮定により
\[
 h^0(E)\otimes_DK=\Omega^1_{D/R}\otimes_DK=0.
\]
従って $F^{-1}\otimes_DK\to F^0\otimes_DK$ は同次元間の全射、すなわち同型である。
その核は自由 $D$-加群の部分加群であり、$D$ が整域なのでtorsionを持たない。
一方、この核の一般階数は零であるから、核そのものが零となる。
よって $h^{-1}(E)=0$ である。

$\mathbb L_{D/R}$ についても同じ議論を適用できる。
有限平坦な局所完全交叉の相対次元は零なので、
その余接複体は同階数の有限自由加群間の射で表示される。
generically etaleという仮定により一般点で同型となるので
$h^{-1}(\mathbb L_{D/R})=0$ である。
従って $\phi$ は両次数のコホモロジーで同型となり、擬同型である。
\end{proof}

\begin{lemma}[最小二項複体の一次障害]\label{lem:minimal-intertwiner-v24}
局所環 $(D,\mathfrak m)$ 上で、正方行列 $J$ の全成分が $\mathfrak m$ に属するとする。
$J^t$ で表される二項複体が、その双対を一つずらした複体と同型ならば、
ある可逆行列 $U,V$ により
\begin{equation}\label{eq:four-intertwiner-v24}
 UJ=J^tV
\end{equation}
と書ける。従って $J$ の一次部分を $J_1$ とすると、
定数行列 $U_0,V_0$ は可逆で、$U_0J_1=J_1^tV_0$ を満たす。
\end{lemma}

\begin{proof}
有界自由複体間の導来同型は鎖ホモトピー同値で表せる。
両複体は最小であり、$D/\mathfrak m$ へのテンソル後の微分は零である。
従って鎖ホモトピー同値の各次数の写像は剰余体上で同型となり、
Nakayamaの補題により $D$ 上でも可逆である。
鎖写像の等式を、転置、逆行列および符号の吸収を行って書けば
式\eqref{eq:four-intertwiner-v24}を得る。
その一次項には $U,V$ の定数項だけが寄与する。
\end{proof}

\subsection{四変数族全体の非critical性}

\begin{theorem}\label{thm:four-noncritical-v24}
式\eqref{eq:four-noncritical-input-v24}の有限平坦族 $D/R$ について
\[
 \mathbb L_{D/R}\not\simeq\mathbb L_{D/R}^\vee[1].
\]
従って $D/R$ は相対対称perfect obstruction theoryを持たない。
また、任意の相対次元の滑らかな形式 $R$-スキーム上で、
相対almost closed一形式の零スキームとして表示することはできず、
特に相対critical locusではない。
これらの結論は、基底と全空間に同じ個数の自明な形式変数を付け加えても成立する。
\end{theorem}

\begin{proof}
余接複体は $[D^2\xrightarrow{J^t}D^2]$ で表され、その全成分は極大イデアルに属する。
仮に自己双対ならば、補題\ref{lem:minimal-intertwiner-v24}により
可逆な定数行列 $U_0,V_0$ が式\eqref{eq:four-intertwiner-v24}の一次項を満たす。
しかし $D\simeq\mathbb C[[X,Y,w,z]]$ では $X,Y,w,z$ の一次項は独立である。
\[
 U_0=\begin{pmatrix}a&b\\c&d\end{pmatrix},\qquad
 V_0=\begin{pmatrix}e&f\\g&h\end{pmatrix}
\]
と書き、四変数の係数を順に比較する。
$w$ の係数から $g=b$, $h=d=0$、
$z$ の係数から $b=g=0$ を得る。
$Y$ の係数から $e=a$, $f=c=0$、
$X$ の係数から $a=e=0$ を得る。
従って $U_0=V_0=0$ となり、可逆性と矛盾する。
これは八つの未知数に対する階数八の有理線形系であり、
独立な八行の小行列式は $16$ に取れる。

相対対称perfect obstruction theoryの非存在は
補題\ref{lem:symmetric-pot-is-cotangent-v24}から従う。
相対almost closed一形式 $\alpha=\sum_{i=1}^ng_i\,dz_i$ の零スキームが
$D/R$ と同型であったとする。
この場合 $R[[z_1,\ldots,z_n]]$ は正則で、商は $R$ 上有限平坦なので、
$(g_1,\ldots,g_n)$ の高さは $n$ である。
従って $g_1,\ldots,g_n$ は正則列をなし、余接複体は
$(\partial g_i/\partial z_j)$ の転置で表される。
almost closed条件 $d_{/R}\alpha\in(g_1,\ldots,g_n)\Omega^2_{/R}$ により、
この行列は $D$ 上で対称となる。
従って余接複体が自己双対となって矛盾する。
これは相対ambientの次元を二に限定しない議論である。
同じ構成はalmost closed一形式が対称obstruction theoryを与える一般的な構成の、
局所完全交叉の場合に当たる\cite[Example 1.19]{BehrendFantechiSymmetric}。

自明な形式変数を加えた場合も $X,Y,w,z$ の一次項の独立性は保たれるので、
同じ係数比較で矛盾を得る。
\end{proof}

\begin{corollary}\label{cor:actual-four-noncritical-v24}
十三変数環 $A_0$ 上の実際の四枚被覆 $\Lambda/A_0$ は、
相対対称perfect obstruction theoryを持たず、
任意の滑らかな相対ambientにおいて相対almost closed一形式の零スキームではない。
\end{corollary}

\begin{proof}
四枚被覆の全形式標準形により、$\Lambda/A_0$ は $D/R$ に
九つの自明な形式変数を加えた族である。
従って定理\ref{thm:four-noncritical-v24}を適用できる。
\end{proof}

\subsection{二方向を除く三変数sliceの明示的potential}

$(p,q)\in\mathbb C^2\setminus\{(0,0)\}$ を固定し、
\begin{align}
 R_{p,q}&=\mathbb C[[u,v,t]],\\
 D_{p,q}&=\frac{R_{p,q}[[X,Y]]}
 {(XY-u,\ Y^2-2X^2-v-ptX-qtY)}
 \label{eq:four-linear-slice-v24}
\end{align}
と置く。これは $(w,z)=(pt,qt)$ による線形sliceである。
$D_{p,q}\simeq\mathbb C[[X,Y,t]]$ で、$R_{p,q}$ 上有限自由階数四となる。

\begin{theorem}[critical sliceの完全分類]\label{thm:four-critical-slices-v24}
式\eqref{eq:four-linear-slice-v24}の族について、次の条件は同値である。
\begin{enumerate}
\item $p^2+2q^2\ne0$。
\item $D_{p,q}/R_{p,q}$ は相対critical locusである。
\item $D_{p,q}/R_{p,q}$ は、ある滑らかな形式的相対ambient上の
相対almost closed一形式の零スキームである。
\item $D_{p,q}/R_{p,q}$ は相対対称perfect obstruction theoryを持つ。
\item $\mathbb L_{D_{p,q}/R_{p,q}}\simeq
\mathbb L_{D_{p,q}/R_{p,q}}^\vee[1]$。
\end{enumerate}
成立する場合、二変数の多項式potentialを
\begin{align}
 W_{p,q}={}&p\left(XY^2-\frac23X^3\right)
 +q\left(\frac13Y^3-2X^2Y\right)\nonumber\\
 &-\frac t2(pX+qY)^2
 +(4qu-pv)X-(2pu+qv)Y
 \label{eq:four-slice-potential-v24}
\end{align}
と取れる。従って射影的なslice方向 $\mathbb P^1$ のうち、
critical性が失われるのは $p^2+2q^2=0$ の二点だけである。
\end{theorem}

\begin{proof}
まず自己双対性の必要条件を求める。
相対Jacobianは
\[
 J_{p,q}=\begin{pmatrix}Y&X\\-4X-pt&2Y-qt\end{pmatrix}
\]
である。補題\ref{lem:minimal-intertwiner-v24}の定数行列について
$X,Y$ の係数だけを比較すると、一般解は
\begin{equation}\label{eq:slice-central-intertwiners-v24}
 U_0=\begin{pmatrix}-4k&h\\2h&k\end{pmatrix},\qquad
 V_0=U_0^t
\end{equation}
となる。実際、この中心の線形系は八未知数に対して階数六である。
$t$ の係数を追加すると、残る条件は
\[
 pk-qh=0
\]
だけである。$(p,q)\ne(0,0)$ なので $(h,k)=c(p,q)$ となり、
\[
 \det U_0=-2c^2(p^2+2q^2)
\]
を得る。よって自己双対性には $p^2+2q^2\ne0$ が必要である。

逆に $p^2+2q^2\ne0$ と仮定し、
\[
 K_{p,q}=\begin{pmatrix}-4q&p\\2p&q\end{pmatrix},\qquad
 \det K_{p,q}=-2(p^2+2q^2)
\]
と置く。式\eqref{eq:four-slice-potential-v24}の直接微分により、
多項式環上の厳密な恒等式
\begin{equation}\label{eq:four-slice-gradient-v24}
 \binom{\partial_XW_{p,q}}{\partial_YW_{p,q}}
 =K_{p,q}
 \binom{XY-u}{Y^2-2X^2-v-ptX-qtY}
\end{equation}
を得る。$K_{p,q}$ が可逆なので、両側の成分が生成するイデアルは等しい。
従ってこれは $D_{p,q}/R_{p,q}$ の相対critical表示である。
critical表示はalmost closed表示を与え、
この表示の対称Hessianは相対対称perfect obstruction theoryを与える。
また $D_{p,q}$ は整域で、
$\det J_{p,q}=4X^2+2Y^2+ptX-qtY$ は非零なのでgenerically etaleである。
補題\ref{lem:symmetric-pot-is-cotangent-v24}により、
対称perfect obstruction theoryは余接複体の自己双対性を含意する。
以上で同値性が従う。
\end{proof}

\begin{remark}[中心potentialの選択と二つの例外]\label{rem:central-pencil-v24}
中心の二つの関係 $(XY,Y^2-2X^2)$ を定数行列で変換して勾配とする
三次斉次potentialは、定数を除き
\[
 W^{(3)}_{h,k}
 =h\left(XY^2-\frac23X^3\right)
 +k\left(\frac13Y^3-2X^2Y\right)
\]
に限られる。その勾配が中心のイデアル全体を生成する条件は
$h^2+2k^2\ne0$ である。
slice方向 $(p,q)$ は、その一次変形が許す中心potentialの方向を
$[h:k]=[p:q]$ に一意に定める。
このため、特定の中心potentialを最初に固定して得る接空間だけを比較する議論では、
すべてのcritical sliceを捉えられない。

$W^{(3)}_{p,q}$ の二変数三次形式としての判別式は
\[
 \operatorname{disc}(W^{(3)}_{p,q})
 =\frac83(p^2+2q^2)^2
\]
である。従って二つの例外は、選択される中心三次形式そのものが退化する方向に一致する。
この計算は、自己双対性に対する行列式条件と独立に確かめられる。
判別式が非零の二変数三次形式は三つの相異なる一次因子を持ち、
複素線形座標変換と非零定数倍により $X^3+XY^2$ と同値である。
従って非例外方向の中心関数は $D_4$ 型である。
\end{remark}

\begin{corollary}[good sliceの関数変形としての最小性]
\label{cor:four-slice-potential-versal-v24}
$p^2+2q^2\ne0$ とする。
式\eqref{eq:four-slice-potential-v24}の三変数族は、
中心関数 $W^{(3)}_{p,q}$ の関数変形について、
相対形式座標変換および基底だけに依存する関数の加算を同一視した
最小versal変形である。
\end{corollary}

\begin{proof}
中心のJacobi代数は $D_0=\mathbb C[[X,Y]]/(XY,Y^2-2X^2)$ である。
定数を除いた基底 $(X,Y,X^2)$ で、三つのパラメータ微分のclassは
\[
 \begin{pmatrix}
 [\partial_u W_{p,q}]&[\partial_v W_{p,q}]&[\partial_t W_{p,q}]
 \end{pmatrix}
 =\begin{pmatrix}
 4q&-p&0\\
 -2p&-q&0\\
 0&0&-(p^2+2q^2)/2
 \end{pmatrix}.
\]
この行列式は $(p^2+2q^2)^2$ であり、非零である。
従って接写像は定数を除くJacobi代数への同型となる。
形式的versalityも直接帰納的に得られる。
任意の基底上の関数変形に対し、基底の極大イデアルの次数ごとに誤差を考える。
そのJacobi代数の非定数部分は三つのパラメータの変更で吸収でき、
Jacobianイデアル内の部分は $X,Y$ の座標変更で吸収できる。
残る定数部分は基底だけの関数の加算で消える。
完備性により、この操作は全形式の同値を与える。
接空間の次元も三なので最小である。
\end{proof}

\subsection{適用範囲と厳密検証}

本節の非critical性は、有限射 $\Lambda\to A_0$ に相対的な性質である。
元の十一点族 $T/B$ は別の有限射であり、その相対almost closed表示と矛盾しない。
また本節は、元の十一点族の相対potentialについて残った
$b_0\in J_M$、または固定した完備非可換局所代数の巡回potentialの問題を決定しない。

\nolinkurl{four_cover_noncritical_v24.py} はPython標準ライブラリの
有理数演算だけで、八未知数の一次障害の階数と小行列式、
中心の二次元解空間、sliceの唯一の追加条件、
式\eqref{eq:four-slice-gradient-v24}の多項式恒等式、
中心三次形式の判別式、および関数変形の接写像の行列式を再構成する。
$p,q$ は記号変数として扱い、有限個の値への代入で恒等式を代用しない。
保存済み証明書との比較が既定動作であり、通常の検証実行ではファイルを変更しない。

\subsection{非criticalなcorank二被覆を持つcritical族}
\label{subsec:critical-source-noncritical-corank-v25}

元の点族と、そのcorank二の軌跡から得られる有限被覆のcritical性は異なる問題である。
この区別は、単に推論が未完成であるという意味にとどまらない。
以下の明示的な例では、長さ十一の有限平坦族が相対potentialを持つ一方で、
そこから得られる四枚被覆は第\ref{sec:four-noncritical-v24}節の非criticalな全versal族になる。

\begin{proposition}[critical族の中の非criticalな四枚被覆]
\label{prop:critical-source-noncritical-corank-v25}
\[
 \widetilde B=\mathbb C[[a,b,c,u,v,w,z]],\qquad
 R=\mathbb C[[u,v,w,z]]
\]
とし、二変数の多項式
\begin{align}
 \mathcal W={}&\frac12X^2Y^2+\frac15(X^5+Y^5)
 -\frac w6X^3+\frac z{12}Y^3-2uXY\nonumber\\
 &+\frac c2X^2+\frac c4Y^2+\frac v4Y^2+aX+bY
 \label{eq:critical-source-potential-v25}
\end{align}
を取る。次の有限族と、そのcorank二の環を考える：
\[
 \widetilde T=
 \frac{\widetilde B[[X,Y]]}{(\mathcal W_X,\mathcal W_Y)},\qquad
 \widetilde\Lambda=
 \frac{\widetilde T}{(\mathcal W_{XX},\mathcal W_{XY},\mathcal W_{YY})}.
\]
このとき以下が成立する。
\begin{enumerate}
\item $\widetilde T$ は七次元正則局所環であり、
$\widetilde B$ 上有限自由階数十一である。
当然ながら $\widetilde T/\widetilde B$ は全次数で相対critical locusである。
\item $\widetilde\Lambda$ は四次元正則局所環であり、$R$ 上有限自由階数四である。
具体的には
\begin{equation}\label{eq:countermodel-four-cover-v25}
 \widetilde\Lambda\simeq
 \frac{R[[X,Y]]}
 {(XY-u,\;Y^2-2X^2-v-wX-zY+4X^3-8Y^3)}.
\end{equation}
\item $\widetilde\Lambda/R$ は、基底と全空間の形式座標変更を許せば、
\[
 \frac{\mathbb C[[u,v,w,z,X,Y]]}
 {(XY-u,\;Y^2-2X^2-v-wX-zY)}
 \quad\text{over }\mathbb C[[u,v,w,z]]
\]
と同型である。従って $\widetilde\Lambda/R$ は相対対称perfect obstruction theoryを持たず、
任意の滑らかな相対ambientで相対critical locusとして表示することもできない。
\end{enumerate}
さらに全ての環に同じ九個の自明な形式変数を加えると、
元の点族の基底は十六次元、corank二被覆の基底は十三次元となる。
有限自由階数はそれぞれ十一、四のままであり、
corank二被覆は実際の $\Lambda/A_0$ と同じ形式同型型を持つ。
\end{proposition}

\begin{proof}
$\mathcal W_X$ と $\mathcal W_Y$ はそれぞれ $a,b$ に係数一で依存するので、
これらを消去すれば
\[
 \widetilde T\simeq\mathbb C[[X,Y,c,u,v,w,z]]
\]
となる。中心のJacobianイデアルは
\[
 I_0=(X^4+XY^2,\;X^2Y+Y^4)\subset\mathbb C[[X,Y]].
\]
$f=X^4+XY^2$, $g=X^2Y+Y^4$ と置くと、
\[
 (1-XY)XY^3=Yf-X^2g,
\]
\[
 (1-XY)Y^6=(Y^2-XY^3+X^3)g-XYf
\]
である。$1-XY$ は単元なので、局所環上で
\[
 I_0=(X^4+XY^2,\;X^2Y+Y^4,\;XY^3,\;Y^6).
\]
右辺の四多項式は、$X>Y$ の辞書式順序に対するGr\"obner基底をなす。
標準単項式は
\[
 1,X,X^2,X^3,\quad Y,XY,\quad Y^2,XY^2,\quad Y^3,Y^4,Y^5
\]
の十一個である。従って中心の長さは十一。
形式的な有限性と、同次元の正則局所環上の有限Cohen--Macaulay加群の自由性から、
$\widetilde T$ は $\widetilde B$ 上有限自由階数十一となる。
なお中心potentialの像は $-Y^5/10\ne0$ であり、中心classも非定数である。

次にHessianの三成分は
\[
 \mathcal W_{XX}=Y^2+4X^3-wX+c,
 \qquad \mathcal W_{XY}=2XY-2u,
\]
\[
 \mathcal W_{YY}=X^2+4Y^3+\frac z2Y+\frac c2+\frac v2
\]
である。$\mathcal W_{XX}=0$ により $c$ を消し、
$\mathcal W_{XX}-2\mathcal W_{YY}$ を取ると
式\eqref{eq:countermodel-four-cover-v25}を得る。
さらに $u,v$ を消去すると
$\widetilde\Lambda\simeq\mathbb C[[X,Y,w,z]]$ となる。

閉ファイバーで $g_1=XY$、
$g_2=Y^2-2X^2+4X^3-8Y^3$ と置く。
\[
 (Y-8Y^2)g_1-Xg_2=2(1-2X)X^3,
\]
\[
 Yg_2+(2X-4X^2)g_1=(1-8Y)Y^3
\]
より $X^3,Y^3\in(g_1,g_2)$ となる。
逆向きの包含も直ちに確認できるので、形式局所環では
\[
 (g_1,g_2)=(XY,Y^2-2X^2).
\]
従って閉ファイバーは基底 $1,X,Y,X^2$ を持つ長さ四の完全交叉であり、
上と同じ有限性・自由性の議論が適用できる。

最後にこの四変数変形のKodaira--Spencer写像を調べる。
$A=\mathbb C[[X,Y]]/(XY,Y^2-2X^2)$ 上での相対Jacobianは
\[
 J_0=
 \begin{pmatrix}
 Y&X\\
 -4X+12X^2&2Y-24Y^2
 \end{pmatrix}.
\]
$A$ の上記基底で $A^2\to A^2$ を八次正方行列として表示すると、その階数は四である。
さらに $u,v,w,z$ の四方向
\[
 (-1,0)^t,\quad (0,-1)^t,\quad (0,-X)^t,\quad(0,-Y)^t
\]
を加えた行列の階数は八となる。
実際、Jacobianの独立な四列とこれらの四列からなる小行列式を $8$ に取れる。
従って $\dim T^1_A=4$ であり、Kodaira--Spencer写像は同型である。
完全交叉の変形の障害消滅と最小versal変形の形式的一意性により、
この四変数族は第\ref{sec:four-versal-v23}節の四変数最小versal族と形式同型となる。
critical性、almost closed表示の存在、相対対称perfect obstruction theoryの存在は
基底と全空間の同型で保たれるので、第\ref{sec:four-noncritical-v24}節の非存在定理が適用できる。
\end{proof}

\paragraph{この例が示す範囲.}
この七変数族を、元の十六変数の族 $T/B$ と同一視しているわけではない。
命題が排除する推論は、
「corank二の軌跡から得られた四枚被覆が非criticalだから、元の有限点族も非criticalである」
という推論である。同じ非criticalな四枚被覆は、明示的な全potentialを持つ
有限平坦な長さ十一の族からも生じる。
元の族の $\beta=0$ を決定するには、corank二被覆への縮約で失われる情報を使う必要がある。

\paragraph{有限証明書.}
\nolinkurl{critical_source_noncritical_corank_v25.py} は、上の微分・消去恒等式、
中心イデアルの飽和と長さ十一、corank中心の長さ四、
Kodaira--Spencer写像の全射性を有理数上で再構成する。
既定動作は同名証明書との読取り専用照合であり、
形式族の有限平坦性とversal性の議論は上の証明に含まれる。


\section{循環 torsion と全 potential の一元所属}\label{sec:obstruction}

\subsection{二つの加群の一般階数と内在的な torsion}

以後、$\operatorname{Tor}_{A_0}(N)$ は整域 $A_0$ 上の加群 $N$ の torsion 部分加群を表す。
導来関手 $\Tor_i^{A_0}$ とは異なる記号の用法である。

\begin{lemma}\label{lem:rank-bound}
$\rk_{A_0}M=\rk_{A_0}Q=4$ である。
\end{lemma}

\begin{proof}
命題\ref{prop:inputs}(c)により両加群は最小五元生成である。
自由階数五なら $\mathfrak m_{A_0}^2$ を法とした長さは $5\cdot14=70$ となるが、
実際の長さは六十七なので自由ではない。
一般階数が五であれば、最小全射 $A_0^5\to N$ の核は階数零の自由加群の部分加群となり、
$A_0$ が整域であるため零である。これは非自由性に反する。
従って $N=M,Q$ の一般階数は四以下である。
一方、定理\ref{thm:psi-vanishing}と定理\ref{thm:fourfold}により
\[
 M\twoheadrightarrow Q\twoheadrightarrow\Omega^2_{T/B}\simeq\Lambda\simeq A_0^4
\]
であるから、両方の階数は四以上である。
\end{proof}

\begin{theorem}[自由部分と循環障害]\label{thm:cyclic-torsion}
ある非零の真のイデアル $J_M\subset A_0$ と元 $b_0\in\mathfrak m_{A_0}^2$ が存在して、
$A_0$-加群として
\begin{align}
 M&\simeq A_0^4\oplus A_0/J_M,\label{eq:M-cyclic-v23}\\
 Q&\simeq A_0^4\oplus A_0/J_Q,\qquad J_Q=J_M+(b_0),\label{eq:Q-cyclic-v23}\\
 H^1(\Omega^\bullet_{T/B})&\simeq A_0/J_Q\label{eq:H1-cyclic-v23}
\end{align}
が成立する。$\beta$ は $M$ の循環部分の $[b_0]$ に対応する。
二つのイデアルはともに
\begin{equation}\label{eq:torsion-linear-parts}
 \frac{J_M+\mathfrak m_{A_0}^2}{\mathfrak m_{A_0}^2}
 =\frac{J_Q+\mathfrak m_{A_0}^2}{\mathfrak m_{A_0}^2}
 =\langle\ell_1,\ell_2,\ell_3\rangle
\end{equation}
を満たし、特に $\operatorname{ht}J_M,\operatorname{ht}J_Q\ge3$ である。
さらに
\[
 \ker(M\to\Lambda)=\operatorname{Tor}_{A_0}M,\qquad
 H^1(\Omega^\bullet_{T/B})=\operatorname{Tor}_{A_0}Q.
\]
このため $J_M,J_Q$ はそれぞれの torsion 部分の annihilator として内在的に定まる。
\end{theorem}

\begin{proof}
$\Lambda$ は $A_0$-自由なので、$M\twoheadrightarrow\Lambda$ は $A_0$ 上で分裂する。
核を $K_M$ とすれば $M\simeq A_0^4\oplus K_M$ であり、
最小生成数が五であることから $K_M$ は非零の循環加群 $A_0/J_M$ となる。
補題\ref{lem:rank-bound}により $K_M$ は階数零、従って torsion である。
逆に $\Lambda$ は torsion を持たないので $M$ の全 torsion は核に入る。
これが $K_M=\operatorname{Tor}_{A_0}M$ を示す。

$\beta$ は $M\to Q$ の核に属するので $K_M$ に入る。
循環生成元を一つ選び、その座標を $b_0\bmod J_M$ とする。
$K_M$ は $B$-部分加群であるため、$B$ の作用は
\[
 B\longrightarrow\End_{A_0}(A_0/J_M)=A_0/J_M
\]
を通る。従って $B\beta=A_0\beta$ である。
式\eqref{eq:Q-cyclic-v23}と $J_Q=J_M+(b_0)$ が従う。
$Q\to\Lambda$ の核は命題\ref{prop:omega2}により
$H^1(\Omega^\bullet_{T/B})$ なので、式\eqref{eq:H1-cyclic-v23}を得る。
同じ階数の議論により、この核は $Q$ の全 torsion に等しい。
また $\beta\in\mathfrak a^2M$ を上の $A_0$-線形直和射影で読むと
$b_0\bmod J_M\in\mathfrak m_{A_0}^2(A_0/J_M)$ となり、
代表 $b_0$ を $\mathfrak m_{A_0}^2$ に取れる。

$A_1=A_0/\mathfrak m_{A_0}^2$ 上で第四 Fitting イデアルを比較する。
$A_1^4\oplus A_1/\overline J_M$ の第四 Fitting イデアルは $\overline J_M$ であり、
命題\ref{prop:inputs}(c)の表示では $(\ell_1,\ell_2,\ell_3)$ である。
$Q$ にも同じことを適用し、式\eqref{eq:torsion-linear-parts}を得る。
三つの線形項は独立であるため、それらを線形項に持つ各イデアルの三元は
$A_0$ の正則パラメータ系の一部をなす。
従って両イデアルの高さは三以上となる。
\end{proof}

\begin{corollary}[Fitting イデアルによる内在的な判定]\label{cor:intrinsic-fitting}
\[
 J_M=\Fitt_4^{A_0}(M),\qquad J_Q=\Fitt_4^{A_0}(Q).
\]
特に
\begin{equation}\label{eq:beta-fitting-equivalence}
 \beta=0
 \quad\Longleftrightarrow\quad b_0\in J_M
 \quad\Longleftrightarrow\quad
 \Fitt_4^{A_0}(M)=\Fitt_4^{A_0}(Q).
\end{equation}
$b_0\bmod J_M$ は循環生成元の変更で単元倍されるが、最後の条件はその選択によらない。
\end{corollary}

\begin{proof}
自由階数四の加群と循環加群の直和の第四 Fitting イデアルは、
循環加群を定義するイデアルである。
定理\ref{thm:cyclic-torsion}と $J_Q=J_M+(b_0)$ を用いればよい。
\end{proof}

% Proof fragment for integration into manuscript III; no standalone PDF.
\subsection{最小関係数が強制する真の高次関係}

以下では $A=A_0$, $\mathfrak m=\mathfrak m_A$ と略記し、
$B=A[[s,u,v]]$ と書く。ここで $s=t_{15}$, $u=t_5-11t_{15}$,
$v=t_{16}$ である。加群の Betti 数は
$\beta_i^R(N)=\dim_\kk\operatorname{Tor}_i^R(\kk,N)$ とする。

\begin{lemma}[二つの基底についての最小関係数]\label{lem:v24-betti-comparison}
本稿の加群 $M$ と $Q$ について
\[
 \beta_1^B(M)=12,\qquad
 \beta_1^B(Q)=
 \begin{cases}12&\beta=0,\\13&\beta\ne0\end{cases}
\]
である。一方、
\[
 \beta_1^A(M)\ge4,\qquad
 \beta_1^A(Q)\ge4,
 \qquad \beta\ne0\Longrightarrow\beta_1^A(Q)\ge5.
\]
\end{lemma}

\begin{proof}
$W_9$ を第一列に取った Schur 表示を $E=[b\ C]$ とし、
$C:B^{12}\to B^3$, $M=\operatorname{coker}C$ とする。
四次まで持ち上がる零定数中心 class の空間は
$E_4=\kk[W_9]$ である。もし $Cc=0$ なら $(0,c)$ は $E$ の全次数の核元であり、
その中心値は $E_4$ に入る。第一座標が零なので中心値は零である。
従って $\ker C\subset\mathfrak m_BB^{12}$ である。
さらに $C(0)=0$ であるから、この表示は第一関係まで最小であり、
$\beta_1^B(M)=12$ を得る。

$\beta\ne0$ と仮定する。$E(a,c)=0$ の中心値は同じく $\kk[W_9]$ に入る。
$a$ が単元なら $b=-Ca^{-1}c$ となり $\beta=[b]=0$ に反する。
従って全ての核元の中心値は零で、$E$ は十三本の最小関係を与える。
$\beta=0$ なら $Q=M$ なので、いずれの場合も表示された Betti 数を得る。

閉 $A$-ファイバー $M_0=M/\mathfrak mM$ は五次元で、基底を
$g_0,g_1,g_2,h_3,h_4$ と取ると、非零の作用は
\[
 sg_0=h_3,\qquad sh_3=h_4,\qquad ug_1=h_4,\qquad vg_2=h_4
\]
だけである。三作用についての Koszul 複体の鎖群次元は
$(5,15,15,5)$、微分階数は $(2,5,4)$ である。従って
\[
 \dim_\kk H_1(s,u,v;M_0)=15-2-5=8.
\]
$\beta\in\mathfrak m^2M$ なので $Q/\mathfrak mQ=M_0$ も成立する。

$N=M,Q$ に対して、$A$ の十三パラメータと $s,u,v$ の Koszul 二重複体を用いる。
そのスペクトル系列は
\[
 E^2_{p,q}=H_p\!\left(s,u,v;
       \operatorname{Tor}_q^A(\kk,N)\right)
 \ \Longrightarrow\ \operatorname{Tor}_{p+q}^B(\kk,N)
\]
である。全次数一の二項について
\[
 \dim_\kk E^2_{1,0}=8,\qquad
 \dim_\kk E^2_{0,1}\le\beta_1^A(N)
\]
だから、フィルトレーションの次元を比較して
\[
 \beta_1^B(N)\le8+\beta_1^A(N)
\]
を得る。これと先の十二または十三という値を合わせればよい。
この議論はスペクトル系列の退化を仮定しない。
\end{proof}

\begin{proposition}[循環 torsion の余次元は少なくとも四]\label{prop:v24-height-four}
循環分解
\[
 M\simeq A^4\oplus A/J_M,\qquad
 Q\simeq A^4\oplus A/J_Q,
 \quad J_Q=J_M+(b_0)
\]
におけるイデアルは
\[
 \mu_A(J_M)\ge4,\qquad\mu_A(J_Q)\ge4,
 \qquad \operatorname{ht}J_M,\operatorname{ht}J_Q\ge4
\]
を満たす。特に、三つの独立一次関係を形式座標変換で消去した後にも、
十変数正則局所環上で非零の高次関係イデアルが必ず残る。
$\beta\ne0$ なら $\mu_A(J_Q)\ge5$ である。
\end{proposition}

\begin{proof}
循環分解より
\[
 \beta_1^A(M)=\beta_1^A(A/J_M)=\mu_A(J_M),\qquad
 \beta_1^A(Q)=\mu_A(J_Q).
\]
したがって最小生成数の主張は補題\ref{lem:v24-betti-comparison}による。

$J=J_M$ または $J_Q$ とする。既知の一次項は独立な三形式
\[
 \ell_1=t_8+9t_9-4t_{10}+5t_{14},\qquad
 \ell_2=t_{11}-t_{14},\qquad \ell_3=t_{13}-t_{14}
\]
である。それぞれを一次項に持つ $g_1,g_2,g_3\in J$ を選ぶ。
これらは $A$ の正則パラメータ系の一部であり、
$A'=A/(g_1,g_2,g_3)$ は十変数の形式冪級数環である。
$\mu_A(J)\ge4$ より $J\ne(g_1,g_2,g_3)$ であるため、
像 $J'=J/(g_1,g_2,g_3)$ は $A'$ の非零の真のイデアルである。
$A'$ は整域なので $\operatorname{ht}_{A'}J'\ge1$、従って
\[
 \operatorname{ht}_A J=3+\operatorname{ht}_{A'}J'\ge4.
\]
全ての一次項は消去済みなので $J'\subset\mathfrak m_{A'}^2$ でもある。
これは高次関係の存在を保証するが、その最低次数や生成元を特定するものではない。
\end{proof}



\subsection{台の次元と射影次元}

\begin{corollary}\label{cor:projective-dimensions}
\[
 \dim\Supp_BQ=13,\qquad \pd_B\Omega^2_{T/B}=3,
\]
\[
 \pd_BM\ge7,\qquad\pd_BQ\ge7,
 \qquad\pd_B\mathscr R_{T/B}\ge5.
\]
さらに
\[
 \dim\Supp_{A_0}\operatorname{Tor}_{A_0}M\le9,\qquad
 \dim\Supp_{A_0}H^1(\Omega^\bullet_{T/B})\le9.
\]
特に相対 de Rham 第一コホモロジーは非零である。
\end{corollary}

\begin{proof}
命題\ref{prop:omega2}により $Q$ と $\Lambda$ の $B$ 上の台は一致し、
有限射は次元を保つので十三次元である。
$\Lambda$ 自身は十三次元正則局所環で、$B$ 上有限である。
局所 cohomology による depth 比較と Auslander--Buchsbaum 公式から
$\pd_B\Lambda=16-13=3$ となる。

有限 $A_0$-加群 $N=M,Q$ に対し、$B/\Ann_BN$ は有限局所 $A_0$-代数である。
実際、その像は有限 $A_0$-加群 $\End_{A_0}N$ の部分加群である。
従って $N$ 上で両極大イデアルの根基は一致し、
$\depth_BN=\depth_{A_0}N$ となる。
命題\ref{prop:v24-height-four}により循環直和因子の次元は高々九なので、
両 depth は高々九であり、$\pd_BN\ge16-9=7$ を得る。

最後に $I/I^2\to T^2$ の conormal 写像は Jacobian で表される。
$T/B$ が generically etale で、$T$ は整域であるため、この写像は単射である。
$\ker D$ から $T$ への射は $\mathscr R_{T/B}$ と同型になる。
実際、$dT$ が零になる class は $P$ へ持ち上げた後、
$I/I^2$ の一意な元を引けば $D$ で零になる。
従って
\[
 0\longrightarrow\mathscr R_{T/B}\longrightarrow B^{33}
  \xrightarrow{D}B^{22}\longrightarrow Q\longrightarrow0
\]
が完全であり、dimension shifting により
$\pd_B\mathscr R_{T/B}=\pd_BQ-2\ge5$ となる。
$A_0/J_Q$ は最小一元生成の非零加群であるから、最後の主張も従う。
\end{proof}

この非零性は、中心ファイバーの naive de Rham 第一コホモロジーが零であることと矛盾しない。
相対 de Rham 複体の各項は全てが $B$-平坦ではなく、
コホモロジーの形成と閉ファイバーへの基底変更は一般に交換できない。
また、これらの加群分解は $A_0$ 上の分解であり、
三つの残る基底変数の作用を保つ $B$ 上の直和分解は主張しない。

\subsection{全相対 potential と、残る一つの所属問題}

\begin{proposition}\label{prop:potential-membership}
中心の四次項が $x^2y^2/2$ で
\[
 (H_x,H_y)=(F_1,F_2)
\]
というイデアルの等式を満たす $H\in B[[x,y]]$ が存在することは、
式\eqref{eq:beta-fitting-equivalence}の条件と同値である。
存在する場合、中心 class を $[W_9]$ として正規化でき、
代数的で収束する相対 potential を選べる。
\end{proposition}

\begin{proof}
このような $H$ の $P=S/I^2$ における class は $\ker D$ に属する。
全四次の必要条件により、定数 class を除いた中心 class は $[W_9]$ の単元倍である。
係数を一に正規化すると、Schur 表示\eqref{eq:Schur-obstruction}の残る方程式は
\[
 b+C_Mv=0
\]
となる。従って $\beta=0$ である。

逆に $\beta=0$ ならこの線形方程式を解き、$\ker D$ における class を得る。
全四次必要条件の同じ一意性から中心 class は正規化した $[W_9]$ である。
任意の持上げ $H\in S$ は $H_x,H_y\in I$ を満たす。
中心の最低次項では
\[
 \binom{H_x}{H_y}
 =\begin{pmatrix}0&1\\1&0\end{pmatrix}
   \binom{F_1}{F_2}+\text{高次項}
\]
となるので、二つの勾配は $I$ の最小生成系であり、生成するイデアルは $I$ に等しい。
系\ref{cor:intrinsic-fitting}が述べた所属条件との同値を与える。

最後に、関係 slice、十九次の単位消去、$b,C_M$ は代数的形式級数環上で定義される。
この環からその完成への写像は忠実平坦なので、完成後に零となる余核の元 $[b]$ は
既に代数的形式級数環上で零である。
従って $b+C_Mv=0$ はその環上で解ける。
有限点族の代数的基底で class を持ち上げれば $H$ も代数的に選べ、
複素数上の代数的形式級数の収束性により収束する解を得る。
\end{proof}

第\ref{sec:serre-vanishing}節で証明した $\psi=0$ は、
この命題の存在条件を全て解いたことにはならない。
得られたのは $T$ 上の自己 Ext 複体の鎖自己双対性と、corank 二の環の正則性である。
相対 potential には $S=B[[x,y]]$ 上で偏微分を積分できる追加条件が必要であり、
それが現在の $b_0\in J_M$ である。
また、固定した完備非可換局所代数が巡回 potential を持つかという問題とも区別する。

\section{検証の構成と未解決問題}\label{sec:verification}

Serre双対性の有限入力を再構成する検証器は \nolinkurl{code/serre_corank_vanishing_v23.py} である。
この検証器は十四関係から左右の行列を作り、実際の $A_*$ の右関係基底変換、
九つの incidence 恒等式、次数零・一の三対、次数一・二の七対、
次数二・三の五対の消去に使う小行列を再構成する。
中心の $F_1,F_2$ も六次 determinant から作り直し、保存された混合関係の全中心係数と比較する。
さらに $U_0J_2=J_2^tV_0$ の十二本の線形条件を八未知数で解き、
階数七、核が $(\text{swap},\text{swap})$ の一方向であることを確認する。
incidence、全係数環上の三次 Ext、中心の Serre 同型については、
三つの独立な既存検証器も読取り専用で再実行する。

\nolinkurl{code/corank2_support_v22.py} は一次 minor と閉ファイバーの長さ四を、
\nolinkurl{code/four_cover_versal_v23.py} は二十次 Jacobian と
Kodaira--Spencer 全射を再計算する。
有限計算から全次数へ移る箇所は、それぞれ
形式陰関数定理、完備 Nakayama、相対 Serre 双対性、最小 versality で明記した。
特に $\psi=0$ は、高次係数を有限個調べた結果ではなく
単元 $U_{12}$ による全環上の恒等式 $U_{12}\psi=0$ の帰結である。

\nolinkurl{code/four_cover_discriminant_v24.py} はtrace判別式、norm、Saito行列を、
\nolinkurl{code/four_cover_noncritical_v24.py} は一次自己双対障害と全記号potentialを再計算する。
前者の乗法行列は、独立なGröbner基底による実装でも再構成した。
\nolinkurl{code/verify_torsion_height_v24.py} は五次元ファイバーのKoszul微分と
一次関係の階数を独立な有理数演算で確認する。
高さ四という結論には、さらに本文のBetti数比較と正則パラメータの消去を用いる。

\nolinkurl{code/critical_source_noncritical_corank_v25.py} は、
命題\ref{prop:critical-source-noncritical-corank-v25}の明示potentialから微分と消去式を作り、
中心の長さ十一、corank中心の長さ四、Kodaira--Spencer写像の同型を有理数上で確認する。
この例の有限平坦性と全形式標準形は、本文の局所環の正則性と最小versalityから従う。

残る存在問題には、選んだ循環生成元での一つの元 $b_0$ と
内在的な理想 $J_M=\Fitt_4^{A_0}(M)$ を用いる。
三つの独立な線形項を持つ $J_M$ の元を正則パラメータに選べば、
その三方向を消去した十変数環上で残るイデアル所属を調べられる。
これは原理的な変数削減であり、その後の十変数環に非零の高次関係イデアルが必ず残る。
従って $J_M$ は三元生成ではなく、$A_0/J_M$ は正則ではない。
その高次生成元の最初の次数と、$b_0$ との関係が次の具体的な検討対象になる。
次の自然な課題は、この所属の全次数証明、または非所属を示す厳密な商・特殊化を得ることである。
有限次数の相対解をさらに一段延ばすだけでは、ここで残った問題は決定されない。

\appendix
\section{代数 $A_*$ の十四関係}\label{app:relations}

積は道を左から右へ読む記法とする。
左関係は次の七本である。
\begin{align*}
L_0={}&x_1u_{33}+x_3u_{13}-x_1u_{11}-2x_2u_{12}+x_2u_{22}-x_3u_{22},\\
L_1={}&x_1u_{22}-x_2u_{12},\\
L_2={}&x_2u_{33}-x_3u_{23}+x_1u_{11}-x_3u_{11}-x_2u_{11}-x_2u_{12},\\
L_3={}&x_2u_{11}+x_1u_{12},\\
L_4={}&x_3u_{22}+x_2u_{23}+x_1u_{11}-x_2u_{11},\\
L_5={}&x_3u_{11}-x_1u_{13}-x_2u_{12}+x_2u_{22},\\
L_6={}&x_1u_{23}+x_2u_{13}+x_3u_{12}.
\end{align*}
右関係は次の七本である。
\begin{align*}
R_0={}&-u_{33}z_1+u_{13}z_3+u_{11}z_1+2u_{11}z_2+u_{22}z_2+u_{22}z_3,\\
R_1={}&-u_{22}z_1-u_{12}z_2,\\
R_2={}&-u_{33}z_2-u_{23}z_3+u_{11}z_1+u_{11}z_3-u_{11}z_2-u_{12}z_2,\\
R_3={}&-u_{11}z_2+u_{12}z_1,\\
R_4={}&-u_{22}z_3+u_{23}z_2+u_{11}z_1+u_{11}z_2,\\
R_5={}&-u_{11}z_3-u_{13}z_1-u_{12}z_2+u_{22}z_2,\\
R_6={}&u_{23}z_1+u_{13}z_2+u_{12}z_3.
\end{align*}
この付録は $A_*$ を定義するためのものであり、十六変数 slice のすべての高次係数を
有限の多項式として置き換えるものではない。
slice 自体は第\ref{sec:input}節の incidence 条件と固定した形式陰関数によって定義される。

\begin{thebibliography}{99}
\bibitem{LO}
K.~Loginov and D.~Orlov,
\emph{Three-dimensional noncommutative Fano minifolds}, arXiv:2610.31250.
\url{https://arxiv.org/abs/2610.31250}.
\bibitem{Shklyarov}
D.~Shklyarov,
\emph{On Serre duality for compact homologically smooth DG algebras},
arXiv:math/0702590 (2007).
\url{https://arxiv.org/abs/math/0702590}.
\bibitem{PandharipandeThomasAlmostClosed}
R.~Pandharipande and R.~P.~Thomas,
\emph{Almost closed 1-forms},
Glasgow Mathematical Journal \textbf{56} (2014), 169--182.
\url{https://doi.org/10.1017/S0017089513000141}.
Preprint: \url{https://arxiv.org/abs/1204.3958}.
\bibitem{IshikawaOpenings}
G.~Ishikawa,
\emph{Openings of differentiable map-germs and unfoldings},
Topics on Real and Complex Singularities, World Scientific (2014), 87--113.
Preprint: Hokkaido University Preprint Series in Mathematics 1016 (2012).
\url{https://doi.org/10.14943/84162}.
\bibitem{IshikawaFrontal}
G.~Ishikawa, \emph{Frontal singularities and related problems},
\S\,1.4, author manuscript.
\url{https://www.math.sci.hokudai.ac.jp/~ishikawa/Ishikawa-frontal-handbook-final.pdf}.
\bibitem{BEvB}
R.-O.~Buchweitz, W.~Ebeling and H.-C.~Graf von Bothmer,
\emph{Low-dimensional Singularities with Free Divisors as Discriminants},
J.~Algebraic Geom. \textbf{18} (2009), 371--406.
\url{https://arxiv.org/abs/math/0612119}.
\bibitem{BehrendFantechiSymmetric}
K.~Behrend and B.~Fantechi,
\emph{Symmetric obstruction theories and Hilbert schemes of points on threefolds},
Algebra \& Number Theory \textbf{2} (2008), 313--345.
\url{https://doi.org/10.2140/ant.2008.2.313}.
\bibitem{Stacks}
The Stacks Project Authors, \emph{The Stacks Project},
Tags 00R4 (miracle flatness), 00DV (Nakayama's lemma),
00MC (faithfully flat completion), 0A1W (algebraic power series and henselization),
08SH (cotangent complex of a local complete intersection).
\url{https://stacks.math.columbia.edu}.
\bibitem{Companion}
本研究の計算補遺, \emph{Veronese--Kronecker deformation bundle}, 第25版 (2026).
四枚被覆の新検証器は \nolinkurl{code/four_cover_discriminant_v24.py} と
\nolinkurl{code/four_cover_noncritical_v24.py}、
高さの補強は \nolinkurl{code/verify_torsion_height_v24.py} による。
明示potentialを持つ比較族は \nolinkurl{code/critical_source_noncritical_corank_v25.py} で検証する。
Serre双対性とversal性の検証器も含む。
既存の関係、入力、独立検証器は \nolinkurl{archive/v21/code/} に収録。
\end{thebibliography}
\end{document}
